% Μαθηματικά Α' Γυμνασίου -- Τεύχος Α
% Μεταγλώττιση: xelatex teyxos_a.tex (δύο φορές, για το περιεχόμενο)
\documentclass[11pt,twoside,openright]{book}
\usepackage{glossa}
\hypersetup{pdftitle={Μαθηματικά Α' Γυμνασίου, Τεύχος Α},pdfauthor={Ιωάννης Ψωμιάδης}}

% ---------- περιεχόμενα ----------
\usepackage{tocloft}
\renewcommand{\cftpartfont}{\bfseries\color{gblue}}
\renewcommand{\cftchapfont}{\normalfont}
\renewcommand{\cftchappagefont}{\normalfont}
\setlength{\cftbeforepartskip}{1.4em}
\setlength{\cftbeforechapskip}{0.25em}
\renewcommand{\cftdot}{.}
\renewcommand{\cftchapleader}{\cftdotfill{\cftdotsep}}
\renewcommand{\cftpartpresnum}{}
\renewcommand{\contentsname}{Περιεχόμενα}
\renewcommand{\cftpartdotsep}{\cftnodots}
\setcounter{tocdepth}{0}
\let\oldpart\part
\newcommand{\epanalipsi}[1]{\clearpage\markboth{Επανάληψη ενότητας #1}{Επανάληψη ενότητας #1}%
  \addcontentsline{toc}{chapter}{\hspace{1.5em}Επανάληψη ενότητας #1}%
  \setcounter{exer}{0}\renewcommand{\exprefix}{Ε#1}\vspace*{14mm}{\color{gblue}\bfseries\large Επανάληψη}\par\vspace{3mm}%
  {\chaptertitlefont Επανάληψη ενότητας #1\par}\vspace{6mm}}

\begin{document}
\frontmatter
\pagestyle{empty}
% ---------- εξώφυλλο ----------
\thispagestyle{empty}
\begin{center}
\vspace*{26mm}
{\color{ggray}\small Η Γλώσσα των Μαθηματικών}\par\vspace{16mm}
{\bfseries\fontsize{24}{30}\selectfont Μαθηματικά Α’ Γυμνασίου\par}\vspace{9mm}
{\small Μέρος Α΄: Αριθμητική – Άλγεβρα \ · \ Μέρος Β΄: Γεωμετρία\par}\vspace{9mm}
{\color{gblue}\large Βοήθημα με νόημα: ασκήσεις που θέλεις να λύσεις\par}\vspace{40mm}
Ιωάννης Ψωμιάδης\par\vspace{2mm}
{\color{ggray}\scriptsize Σύμφωνα με το Σχολικό Βιβλίο της Α’ Γυμνασίου\par}\vspace{6mm}
{\small Οκτώβριος 2026\par}
\end{center}
\cleardoublepage
\pagestyle{fancy}
\fancyhead[RO]{\color{ggray}\small Πριν ξεκινήσεις}
\fancyhead[LE]{\color{ggray}\small Η Γλώσσα των Μαθηματικών}
\pagenumbering{roman}

\chapter*{}\vspace{-22mm}
{\color{gblue}\bfseries\Large Πριν ξεκινήσεις}\par\vspace{4mm}
{\color{gblue}\bfseries\large Ένα γράμμα σε σένα}\par

Μπήκες στο γυμνάσιο. Άλλο κτίριο, άλλα πρόσωπα, ένα πρόγραμμα με διαφορετικό μάθημα σε κάθε ώρα και έναν καθηγητή για καθένα. Αν έχεις διαβάσει τα βιβλία της Ε’ και της ΣΤ’, θα αναγνωρίσεις τα παιδιά, τον τρόπο και τη φιλοσοφία. Αν όχι, δεν πειράζει: όλα όσα χρειάζεσαι ξαναλέγονται εδώ, όπου χρειάζεται.

Τα περισσότερα βιβλία μαθηματικών σου λένε πώς λύνεται μια άσκηση. Αυτό το βιβλίο θέλει να σου πει κάτι άλλο: γιατί αξίζει να τη λύσεις.

Ξέρεις ήδη να μιλάς ελληνικά. Ίσως λίγο αγγλικά. Με αυτές τις γλώσσες περιγράφεις τον κόσμο: τι βλέπεις, τι θέλεις, τι σε ενοχλεί, τι σε κάνει να γελάς. Τα μαθηματικά είναι κι αυτά μια γλώσσα. Περιγράφουν τον ίδιο κόσμο, απλώς με άλλες λέξεις. Όταν γράφεις $5>3$, λες την ίδια πρόταση με το «το πέντε είναι μεγαλύτερο από το τρία». Το σύμβολο $>$ είναι μια λέξη. Το $=$ είναι μια λέξη. Ακόμα και ο αριθμός 1.000.000 είναι μια λέξη, που στα ελληνικά τη λέμε «ένα εκατομμύριο».

Στο δημοτικό έμαθες λέξεις. Στην ΣΤ’ το λεξιλόγιό σου μεγάλωσε. Φέτος μαθαίνεις γραμματική: τους κανόνες που λένε ποιες προτάσεις επιτρέπονται και ποιες όχι. Γιατί το $-3$ υπάρχει, αφού δεν μπορείς να έχεις «μείον τρία μήλα»; Γιατί $(-1)\cdot(-1)=1$; Γιατί ισχύει το «κάθε τρίγωνο έχει άθροισμα γωνιών 180°» και πώς το ξέρουμε χωρίς να τα μετρήσουμε όλα; Φέτος για πρώτη φορά δεν θα σε ρωτήσω μόνο «πόσο κάνει;». Θα σε ρωτήσω και «γιατί;».

Σε αυτό το βιβλίο θα μάθεις τη γλώσσα των μαθηματικών όπως μαθαίνει κανείς κάθε γλώσσα: λέξη λέξη, κανόνα κανόνα, με παραδείγματα από τη ζωή σου και με πολλή εξάσκηση. Δεν θα σου ζητήσω να αποστηθίσεις κανόνες που δεν καταλαβαίνεις. Κάθε κανόνας θα έχει έναν λόγο ύπαρξης, και θα τον δεις πριν δεις τον κανόνα.

Ένα ακόμη. Οι ασκήσεις εδώ δεν σου ζητούν να λύσεις κάτι επειδή έτσι πρέπει. Σου ζητούν να βοηθήσεις κάποιον, να κερδίσεις ένα στοίχημα, να αποδείξεις ότι έχεις δίκιο, να λύσεις μια διαφωνία. Τα μαθηματικά είναι το εργαλείο. Ο λόγος είναι πάντα δικός σου.

{\color{gblue}\bfseries\large Η παρέα}\par

Το Σεπτέμβρη η παρέα της ΣΤ’ βρέθηκε στο γυμνάσιο, στο ίδιο τμήμα, την Α3. Μαζί της βρέθηκαν και δύο παιδιά από το άλλο δημοτικό της περιοχής, και ένα κορίτσι που δεν ήξερε κανέναν. Οι δύο γειτονιές έγιναν μία τάξη. Η παρέα αυτή θα σε συνοδεύει σε όλο το γυμνάσιο, και τα προβλήματά της θα μεγαλώνουν μαζί της.

\begin{itemize}[leftmargin=1.4em,itemsep=0.2ex]
\item Η \textbf{Ζωή} ρωτάει πάντα «γιατί;» και δεν σταματάει στην πρώτη απάντηση.
\item Ο \textbf{Ηλίας} μετράει, δοκιμάζει, φτιάχνει πίνακες. Όταν κάνει λάθος, το βρίσκει μόνος του, αρκεί να τον αφήσεις.
\item Η \textbf{Μυρτώ} έχει το πρακτικό μυαλό της κουζίνας της γιαγιάς: αν μια απάντηση δεν βγάζει νόημα στην πραγματικότητα, δεν τη δέχεται.
\item Ο \textbf{Στέλιος} μαθαίνει έναν κανόνα και τον εφαρμόζει με ενθουσιασμό παντού, ακόμα και εκεί που δεν ταιριάζει.
\item Η \textbf{Φανή} κρατάει σημειώσεις για όλους και ρωτάει όσα οι άλλοι δεν τολμούν.
\item Ο \textbf{Βασίλης}, ο πρώην πρόεδρος του μαθητικού συμβουλίου, έχει πάντα άποψη και την υποστηρίζει με σιγουριά, κάποτε και με περισσότερη απ' όση χρειάζεται.
\item Ο \textbf{Νίκος}, από το άλλο δημοτικό, παίζει μπάσκετ και βιάζεται. Ό,τι γίνεται γρήγορα, το κάνει πρώτος. Ό,τι γίνεται σωστά, το κάνει δεύτερος.
\item Η \textbf{Αγγελική}, από το άλλο δημοτικό, ζωγραφίζει και παίζει μουσική. Βλέπει σχήματα και μοτίβα εκεί που οι άλλοι βλέπουν μόνο αριθμούς.
\item Η \textbf{Πολυξένη} είναι η καινούργια. Ήρθε το καλοκαίρι από τη Θεσσαλονίκη. Κρατάει ένα τετράδιο με κωδικούς και μοτίβα, μιλάει λίγο και ρωτάει σπάνια. Όταν ρωτάει, η ερώτησή της είναι αυτή που χρειαζόταν να ακουστεί.
\item Η \textbf{γιαγιά Ρήνη}, ο \textbf{παππούς} και ο \textbf{θείος Λευτέρης} είναι πάντα εκεί, με τον κήπο, το χωριό και τα πλακάκια τους.
\end{itemize}

Και υπάρχει η \textbf{κυρία Αλεξάνδρα}, η καθηγήτρια των μαθηματικών τους. Δεν δίνει ποτέ έτοιμη απάντηση. Δίνει ερώτηση.

\clearpage
{\color{gblue}\bfseries\large Πώς είναι φτιαγμένο το βιβλίο}\par

Όλες οι παράγραφοι έχουν την ίδια σειρά. Μόλις τη μάθεις, θα ξέρεις πάντα πού βρίσκεσαι.

\begin{enumerate}[leftmargin=2em,itemsep=0.5ex]
\item \textbf{Η ανάγκη.} Μια μικρή ιστορία με ένα πρόβλημα που δεν λύνεται χωρίς το εργαλείο του κεφαλαίου. Πρώτα νιώθεις ότι κάτι λείπει. Μετά σου το δίνω.
\item \textbf{Η λέξη.} Η νέα έννοια στα ελληνικά: τι σημαίνει, πού τη συναντάς. Μετά ο ακριβής ορισμός της.
\item \textbf{Η μετάφραση.} Ένας πίνακας με δύο στήλες. Αριστερά μια πρόταση στα ελληνικά, δεξιά η ίδια πρόταση στα μαθηματικά.
\item \textbf{Η θεωρία.} Οι κανόνες, με όλη την ακρίβεια που χρειάζονται. Είναι γραμμένοι σύμφωνα με το Σχολικό Βιβλίο σου, ώστε να βρίσκεις αμέσως την αντίστοιχη παράγραφο στην τάξη.
\item \textbf{Λυμένα παραδείγματα.} Κάθε λύση γράφεται ως σκέψη, σε τρία βήματα: Τι ξέρω. Τι θέλω. Πώς τα συνδέω. Στο τέλος ελέγχουμε αν η απάντηση βγάζει νόημα.
\item \textbf{Ασκήσεις σε τρεις ζώνες.}
  \begin{itemize}[leftmargin=1.2em,itemsep=0.2ex]
  \item \emph{Το κατάλαβα:} εφαρμόζεις τον κανόνα κατευθείαν.
  \item \emph{Το χρησιμοποιώ:} λύνεις ένα πρόβλημα από τη ζωή.
  \item \emph{Το σκέφτομαι:} μια ερώτηση που σε βάζει να σκεφτείς πιο μακριά, συχνά προς το επόμενο κεφάλαιο.
  \end{itemize}
\item \textbf{Λύσεις.} Όλες στο τέλος του βιβλίου, γραμμένες με τον ίδιο τρόπο σκέψης. Κοίτα τες αφού προσπαθήσεις, όχι πριν.
\end{enumerate}

Το βιβλίο χωρίζεται σε δύο μέρη, όπως και το Σχολικό σου: \textbf{Μέρος Α΄, Αριθμητική – Άλγεβρα}, και \textbf{Μέρος Β΄, Γεωμετρία}. Κάθε μέρος έχει κεφάλαια και κάθε κεφάλαιο \emph{παραγράφους} με την ίδια αρίθμηση που έχουν στο Σχολικό: η παράγραφος 1.4 του βιβλίου αυτού αντιστοιχεί στην Α.1.4 του Σχολικού. Οι ασκήσεις αριθμούνται με την παράγραφο: η άσκηση 1.4.7 είναι η έβδομη της παραγράφου 1.4.

Στο τέλος κάθε κεφαλαίου θα βρεις μια \textbf{Επανάληψη} με έναν «χάρτη»: ένα πίνακα που σου λέει σε ποιες παραγράφους να γυρίσεις, αν κάποια άσκηση σε δυσκόλεψε.

{\color{gblue}\bfseries\large Για τον γονέα}\par

Το βιβλίο αυτό γράφτηκε ώστε να το διαβάζετε μαζί με το παιδί σας. Δεν χρειάζεται να θυμάστε τα μαθηματικά του γυμνασίου.

Ξεκινήστε από την «Ανάγκη» και τη «Λέξη». Είναι γραμμένες σε απλά ελληνικά. Αν τις καταλάβετε εσείς, μπορείτε να τις εξηγήσετε. Η «Μετάφραση» δείχνει πώς η ίδια ιδέα γράφεται με σύμβολα, και εκεί συνήθως κολλάνε τα παιδιά: όχι στην ιδέα, αλλά στη μετάφρασή της.

Τρεις συμβουλές από δέκα χρόνια διδασκαλίας:

\begin{itemize}[leftmargin=1.4em,itemsep=0.4ex]
\item Ζητήστε από το παιδί να σας εξηγήσει μια λύση με δικά του λόγια. Αν μπορεί να την πει στα ελληνικά, την έχει καταλάβει.
\item Ένα λάθος είναι πληροφορία, όχι αποτυχία. Ρωτήστε «πώς το σκέφτηκες;» πριν πείτε «είναι λάθος».
\item Οι ασκήσεις της ζώνης «Το σκέφτομαι» δεν είναι υποχρεωτικές. Είναι για τις μέρες που υπάρχει όρεξη.
\end{itemize}

{\color{gblue}\bfseries\large Ποιος το έγραψε}\par

Σπούδασα μαθηματικά και δίδαξα ως φροντιστής για περίπου δέκα χρόνια. Τα επόμενα χρόνια δούλεψα στη βιομηχανία, κάνοντας μη καταστροφικούς ελέγχους: ελέγχω σωλήνες, δοχεία, δεξαμενές και συγκολλήσεις χωρίς να τα κόψουμε, για να ξέρουμε αν είναι ασφαλή.

Σε αυτά τα χρόνια έμαθα κάτι που δεν μου το είχε πει κανένα σχολικό βιβλίο. Τα μαθηματικά δεν είναι μάθημα. Είναι ο τρόπος που οι άνθρωποι συμφωνούν για το τι συμβαίνει. Όταν δύο μηχανικοί διαφωνούν, δεν κερδίζει αυτός που φωνάζει πιο δυνατά. Κερδίζει αυτός που φέρνει τη μέτρηση.

Αυτό το βιβλίο είναι η προσπάθειά μου να σου δώσω νωρίς αυτό το εργαλείο. Και, αν τα καταφέρω, λίγη από την αγάπη που έχω για αυτή τη γλώσσα.

\hfill Ιωάννης Ψωμιάδης

\cleardoublepage
\fancyhead[RO]{\color{ggray}\small Περιεχόμενα}
{\color{gblue}\bfseries\Large Περιεχόμενα}\par\vspace{3mm}
\makeatletter
\@starttoc{toc}
\makeatother


\mainmatter
\pagestyle{fancy}
\fancyhead[RO]{\color{ggray}\small\leftmark}
\enotita{1}{Οι φυσικοί αριθμοί}{Από το λεξιλόγιο στη γραμματική}
\kefalaio{Φυσικοί αριθμοί: η σειρά που δεν τελειώνει}{Κεφάλαιο Α.1.1, «Φυσικοί αριθμοί – Διάταξη φυσικών – Στρογγυλοποίηση»}

\begin{anagke}
Δεύτερη μέρα στο γυμνάσιο. Στην είσοδο μιας αίθουσας ένα χαρτί λέει: «Το σχολείο μας έχει 387 μαθητές». Ο Άρης το φωτογραφίζει και στέλνει στον αδελφό του τον Θοδωρή: «Το σχολείο έχει 400 παιδιά!». Ο Θοδωρής απαντάει με ένα χαμογελαστό «στρογγύλεψες».

Στο μάθημα η κυρία Αλεξάνδρα δεν ανοίγει βιβλίο. Γράφει στον πίνακα τρεις αριθμούς: τους μαθητές του σχολείου, 387, το υψόμετρο του Ολύμπου, 2.917 μέτρα, και τον αριθμό του τηλεφώνου της γραμματείας. «Σας ζητάω να τους πείτε σε ένα φίλο σας. Ποιον θα πείτε ακριβώς και ποιον περίπου;»

«Όλους ακριβώς», λέει ο Στέλιος. «Τι νόημα έχει το “περίπου”;». Η Ζωή σηκώνει το χέρι. «Αν στον φίλο σου πεις ότι περπάτησες 8.412 βήματα, θα θυμάται μόνο το 8 και θα μπερδεύεται με τα υπόλοιπα. Το “περίπου 8.000” το κρατάς». «Και πώς ξέρω ότι όλοι θα βγάλουμε το ίδιο “περίπου”;» ρωτάει η Λίνα. «Αν ο Άρης κάνει το 387 τετρακόσια και εγώ τριακόσια ογδόντα;»

Η κυρία Αλεξάνδρα χαμογελάει. Η ερώτηση της Λίνας είναι ακριβώς αυτή που ήθελε να ακούσει. Πότε επιτρέπεται το «περίπου», και πώς γίνεται με κανόνα, ώστε όλοι να κάνουμε το ίδιο;
\end{anagke}

\begin{lexi}{φυσικός, στρογγυλοποίηση}{φυσικός, στρογγυλοποίηση}
Οι \lexiword{φυσικοί} αριθμοί λέγονται έτσι επειδή είναι οι πρώτοι που βρίσκει κανείς στη φύση, όταν μετράει πράγματα: μήλα, παιδιά, βήματα. Κανείς δεν μέτρησε ποτέ «ενάμισι παιδί». Το μηδέν μπήκε στην παρέα αργότερα. Στο σχολείο θα το βρεις μέσα στους φυσικούς.

Η \lexiword{στρογγυλοποίηση} είναι το να κάνεις έναν αριθμό «στρογγυλό», δηλαδή να τον αλλάξεις με έναν κοντινό του που τελειώνει σε μηδενικά: 387 γίνεται 400. Στρογγυλός σημαίνει τακτοποιημένος, χωρίς αιχμές. Δεν είναι μαγικό: είναι \emph{συμφωνία}. Όλοι κάνουμε τα ίδια βήματα και βγάζουμε το ίδιο αποτέλεσμα.

Δύο ακόμα λέξεις που θα ξανασυναντήσεις: \lexiword{άρτιος} είναι ο αριθμός που «ταιριάζει» τέλεια σε ζευγάρια, και \lexiword{περιττός} αυτός που «περισσεύει» ένας. Τους λέμε και ζυγούς και μονούς.
\end{lexi}

\begin{metafrasi}
Ο αριθμός που ακολουθεί το 9. & $9+1=10$\\
Το πέντε είναι μεγαλύτερο από το τρία. & $5>3$\\
Ο αριθμός $x$ βρίσκεται ανάμεσα στο 3 και στο 8. & $3<x<8$\\
Τρία εκατομμύρια, σαράντα χιλιάδες και επτά. & 3.040.007\\
Το 9.573.842 στρογγυλοποιημένο στις εκατοντάδες. & $9.573.842\approx 9.573.800$\\
Ο αριθμός είναι άρτιος. & Διαιρείται με το 2\\
\end{metafrasi}

\tmima{Η θεωρία}

\begin{orismos}{φυσικοί αριθμοί}
Οι αριθμοί $0,\,1,\,2,\,3,\,4,\,\dots,\,99,\,100,\,\dots$ λέγονται \textbf{φυσικοί αριθμοί}. Κάθε φυσικός έχει έναν επόμενο και, εκτός από το 0, έναν προηγούμενο.
\end{orismos}

Οι φυσικοί χωρίζονται σε δύο κατηγορίες: \textbf{άρτιους} (ή ζυγούς), που διαιρούνται με το 2, και \textbf{περιττούς} (ή μονούς), που δεν διαιρούνται. Ο άρτιος τελειώνει πάντα σε 0, 2, 4, 6 ή 8.

\tmimasmall{Δέκα ψηφία, άπειροι αριθμοί}

Με τα δέκα ψηφία $0,\,1,2,\dots,9$ γράφουμε όσους αριθμούς θέλουμε. Αυτό γίνεται επειδή η αξία ενός ψηφίου εξαρτάται και από τη \textbf{θέση} του, δηλαδή από τη \textbf{δεκαδική τάξη}: μονάδες, δεκάδες, εκατοντάδες, χιλιάδες, δεκάδες χιλιάδων, εκατοντάδες χιλιάδων, εκατομμύρια. Το 0 κρατάει μια θέση άδεια: στο 5.780.901 τα δύο μηδενικά λένε ότι δεν υπάρχουν ούτε χιλιάδες ούτε δεκάδες.

\tmimasmall{Διάταξη: ποιος είναι μεγαλύτερος}

Γράφουμε $=$ για «ίσος με», $<$ για «μικρότερος από» και $>$ για «μεγαλύτερος από». Δύο φυσικούς μπορούμε πάντα να τους συγκρίνουμε.

\begin{kanonas}{πώς συγκρίνω δύο φυσικούς}
\begin{enumerate}[leftmargin=1.6em,itemsep=0.2ex,topsep=2pt]
\item Μετράω τα ψηφία. Ο αριθμός με περισσότερα ψηφία είναι μεγαλύτερος (αν δεν ξεκινάει με 0).
\item Αν έχουν ίδιο πλήθος ψηφίων, τα κοιτάζω από αριστερά προς τα δεξιά. Στο πρώτο ψηφίο που διαφέρουν, μεγαλύτερος είναι αυτός με το μεγαλύτερο ψηφίο.
\end{enumerate}
\end{kanonas}

Έτσι $90.876>86.945$, γιατί στο πρώτο ψηφίο το 9 είναι μεγαλύτερο από το 8, και $345<5.690$, γιατί το 345 έχει λιγότερα ψηφία. Τους φυσικούς μπορούμε να τους βάλουμε σε σειρά: \textbf{αύξουσα}, από τον μικρότερο στον μεγαλύτερο, ή \textbf{φθίνουσα}, αντίστροφα. Δεν υπάρχει μεγαλύτερος φυσικός αριθμός, γιατί κάθε φυσικός έχει επόμενο. Ο μικρότερος είναι το 0.

\tmimasmall{Η αριθμητική ευθεία}

Η διάταξη μας επιτρέπει να βάλουμε τους φυσικούς πάνω σε μια ευθεία. Διαλέγουμε ένα σημείο Ο για το 0 (την αρχή) και ένα σημείο Α δεξιά του για το 1. Το τμήμα ΟΑ είναι η \textbf{μονάδα μέτρησης}. Με το ίδιο μήκος βρίσκουμε τα σημεία των 2, 3, 4, \dots

\begin{center}
\begin{tikzpicture}[x=1.15cm,y=1cm]
\draw[-{Stealth}] (-0.3,0)--(9.6,0);
\foreach \i in {0,...,8}{\draw (\i,-0.12)--(\i,0.12); \node[below,font=\small] at (\i,-0.12) {\i};}
\node[above,font=\small\bfseries,gblue] at (0,0.15) {Ο};
\node[above,font=\small\bfseries,gblue] at (1,0.15) {Α};
\draw[<->,gblue] (0,0.75)--(1,0.75) node[midway,above,font=\scriptsize] {μονάδα};
\end{tikzpicture}
\end{center}

\tmimasmall{Στρογγυλοποίηση}

\begin{kanonas}{στρογγυλοποίηση φυσικού}
\begin{enumerate}[leftmargin=1.6em,itemsep=0.2ex,topsep=2pt]
\item Βρίσκω την τάξη στην οποία θα στρογγυλοποιήσω (δεκάδες, εκατοντάδες, χιλιάδες, \dots).
\item Κοιτάζω το ψηφίο της \textbf{αμέσως μικρότερης} τάξης.
\item Αν είναι $0,\,1,2,\,3$ ή $4$, το ψηφίο της τάξης μου μένει ίδιο. Αν είναι $5,\,6,7,\,8$ ή $9$, αυξάνεται κατά 1.
\item Όλα τα ψηφία δεξιά από την τάξη μου γίνονται 0.
\end{enumerate}
\end{kanonas}

Το αποτέλεσμα λέγεται \textbf{προσέγγιση} του αριθμού και γράφεται με το σύμβολο $\approx$ (διαβάζεται «περίπου ίσο»). Για παράδειγμα, $2.917\approx 3.000$ στις χιλιάδες και $2.917\approx 2.900$ στις εκατοντάδες.

Η στρογγυλοποίηση ταιριάζει εκεί που μας ενδιαφέρει η τάξη μεγέθους, όπως το πλήθος των ανθρώπων σε μια πόλη. Δεν ταιριάζει εκεί που χρειάζεται ακρίβεια: ο αριθμός του τηλεφώνου, το ΑΦΜ, ο κωδικός ενός λουκέτου.

\begin{prosoxi}
\begin{itemize}[leftmargin=1.2em,itemsep=0.3ex,topsep=2pt]
\item Στρογγυλοποιούμε \textbf{μία φορά}, από τον αρχικό αριθμό. Το 1.449 στις εκατοντάδες είναι 1.400, γιατί το ψηφίο των δεκάδων είναι 4. Αν το κάνεις σε δύο βήματα, πρώτα στις δεκάδες (1.450) και μετά στις εκατοντάδες, βγάζεις 1.500, που είναι λάθος.
\item Όταν το ψηφίο που αυξάνεται είναι 9, η αύξηση «κουβαλάει»: $9.995\approx 10.000$ στις δεκάδες.
\item Όταν μετράς \textbf{από ως}, και τα δύο άκρα μετράνε. Από τη σελίδα 12 ως τη σελίδα 35 είναι $35-12+1=24$ σελίδες, όχι 23.
\end{itemize}
\end{prosoxi}

\tmima{Λυμένα παραδείγματα}

\paradeigma{Οι κωδικοί του ντουλαπιού}
\tixero Το ντουλάπι του Ηλία ανοίγει με έναν τριψήφιο κωδικό. Θυμάται ότι έχει τα ψηφία 4, 0 και 7, καθένα μία φορά, αλλά όχι με ποια σειρά.\par
\tithelo Όλους τους πιθανούς κωδικούς, σε αύξουσα σειρά.\par
\possindeo Δοκιμάζω όλες τις σειρές των τριών ψηφίων: 407, 470, 047, 074, 704, 740. Ο 047 και ο 074 δεν είναι τριψήφιοι αριθμοί, γιατί το 0 μπροστά δεν μετράει. Τους σβήνω.\par
\apantisi 407, 470, 704, 740.\par
\elegxos Είναι τέσσερις, και κανένας δεν ξεκινάει από 0. Ο μικρότερος είναι ο 407, γιατί ξεκινάει από το μικρότερο μη μηδενικό ψηφίο, το 4, και μετά βάζει το μικρότερο που μένει, το 0.

\paradeigma{Στρογγυλοποιώ το ίδιο νούμερο τρεις φορές}
\tixero Ο αριθμός 9.573.842.\par
\tithelo Να τον στρογγυλοποιήσω (α) στις εκατοντάδες, (β) στις χιλιάδες, (γ) στα εκατομμύρια.\par
\possindeo
(α) Τάξη: εκατοντάδες. Η αμέσως μικρότερη τάξη είναι οι δεκάδες, με ψηφίο 4. Επειδή $4<5$, το 8 μένει ίδιο και τα δεξιά γίνονται 0: $9.573.842\approx 9.573.800$.\par
(β) Τάξη: χιλιάδες, ψηφίο 3. Η αμέσως μικρότερη τάξη, οι εκατοντάδες, έχει ψηφίο 8. Επειδή $8\ge 5$, το 3 γίνεται 4: $9.573.842\approx 9.574.000$.\par
(γ) Τάξη: εκατομμύρια, ψηφίο 9. Η εκατοντάδα χιλιάδων έχει ψηφίο 5. Επειδή $5\ge 5$, το 9 αυξάνεται και γίνεται 10: $9.573.842\approx 10.000.000$.\par
\apantisi 9.573.800, 9.574.000 και 10.000.000.\par
\elegxos Κάθε προσέγγιση απέχει λιγότερο από μισή τάξη από τον αρχικό. Για παράδειγμα, στο (β), το 9.573.842 απέχει 158 από το 9.574.000 και 842 από το 9.573.000. Το πρώτο είναι πιο κοντά.

\paradeigma{Πόσες σελίδες;}
\tixero Η Φανή διαβάζει από τη σελίδα 47 ως τη σελίδα 112 του βιβλίου της, και τις δύο μαζί.\par
\tithelo Πόσες σελίδες διαβάζει.\par
\possindeo Αν έβγαζα $112-47=65$, θα είχα αφήσει έξω την πρώτη σελίδα. Για να μετρήσω και τα δύο άκρα, προσθέτω 1: $112-47+1=66$.\par
\apantisi 66 σελίδες.\par
\elegxos Δοκιμή με μικρούς αριθμούς: από τη σελίδα 3 ως τη 5 είναι τρεις σελίδες (3, 4, 5). Ο τύπος δίνει $5-3+1=3$. Σωστά.

\tmima{Ασκήσεις}

\zoni{Το κατάλαβα}{εφαρμογή του κανόνα}

\ex Γράψε με ψηφία: α) διακόσια πέντε \quad β) επτακόσια τριάντα δύο \quad γ) είκοσι χιλιάδες οκτακόσια δεκατρία \quad δ) τέσσερα εκατομμύρια πεντακόσιες χιλιάδες δεκαπέντε.

\ex Γράψε με λέξεις: α) 38.951 \quad β) 5.000.812 \quad γ) 120.003.

\ex Ποιοι είναι οι τρεις προηγούμενοι αριθμοί του 289 και ποιοι οι δύο επόμενοι; Ποιος είναι ο προηγούμενος του 1.000.000 και ποιος ο προηγούμενος του 0;

\ex Γράψε σε αύξουσα σειρά τους αριθμούς 3.515, 4.800, 3.620, 3.508 και 4.801. Μετά γράψε τους ίδιους σε φθίνουσα.

\ex Βάλε το σωστό σύμβολο ($<$, $=$ ή $>$): α) $45\ \square\ 45$ \quad β) $38\ \square\ 36$ \quad γ) $456\ \square\ 465$ \quad δ) $8.765\ \square\ 8.970$ \quad ε) $90.876\ \square\ 86.945$ \quad στ) $345\ \square\ 5.690$.

\ex Σχεδίασε μια αριθμητική ευθεία με αρχή το Ο και μονάδα $ΟΑ=2$ εκ. Τοποθέτησε τα σημεία Β, Γ, Δ, Ε σε αποστάσεις 6, 10, 12 και 14 εκ. από το Ο. Ποιοι αριθμοί αντιστοιχούν στα σημεία αυτά;

\ex Στρογγυλοποίησε στην πλησιέστερη εκατοντάδα: 345, 761, 659, 2.567, 9.532, 31.549 και 8.950.

\ex Στρογγυλοποίησε τον αριθμό 7.568.349 (α) στις δεκάδες, (β) στις εκατοντάδες, (γ) στις χιλιάδες, (δ) στις δεκάδες χιλιάδων, (ε) στα εκατομμύρια.

\ex Σωστό ή λάθος;\par
α) Στον αριθμό 5.780.901 τα δύο μηδενικά δείχνουν ότι δεν υπάρχουν χιλιάδες και δεκάδες.\quad β) Δέκα χιλιάδες είναι μία δεκάδα χιλιάδων.\quad γ) Από τον 32 ως τον 122 υπάρχουν 91 αριθμοί.\quad δ) Από τη σελίδα 12 ως τη σελίδα 35 είναι 24 σελίδες.\quad ε) Σε οκτώ ημέρες από σήμερα, που είναι Πέμπτη, θα είναι Παρασκευή.\quad στ) Ανάμεσα στο 2 και στο 3 δεν υπάρχει άλλος φυσικός αριθμός.\quad ζ) Το 1.449 στρογγυλοποιημένο στις εκατοντάδες είναι 1.500.

\zoni{Το χρησιμοποιώ}{προβλήματα από τη ζωή}

\ex \textbf{Το ντουλάπι της Ζωής.} Ο κωδικός του ντουλαπιού της Ζωής είναι τριψήφιος αριθμός με τα ψηφία 3, 0 και 6, καθένα μία φορά. Γράψε όλους τους δυνατούς κωδικούς από τον μικρότερο στον μεγαλύτερο. Γιατί δεν υπάρχει ο κωδικός 036;

\ex \textbf{Ο Όλυμπος.} Η κορυφή του Ολύμπου, ο Μύτικας, έχει υψόμετρο 2.917 μέτρα. Η Λίνα γράφει σε μια πινακίδα για το σχολικό ημερολόγιο «περίπου 3.000 μ.» και ο Ηλίας «περίπου 2.900 μ.». Ποιος στρογγυλοποίησε σε χιλιάδες και ποιος σε εκατοντάδες; Ποιος από τους δύο λέει κάτι πιο κοντά στην αλήθεια;

\ex \textbf{Τα βήματα του Άρη.} Η εφαρμογή στο κινητό του Άρη γράφει ότι σήμερα έκανε 9.987 βήματα. Ο Άρης λέει «Στρογγυλά, 10.000! Πέτυχα τον στόχο μου». Ο στόχος είναι τουλάχιστον 10.000 βήματα. Έφτασε τον στόχο; Πόσα βήματα του λείπουν; Γιατί εδώ η στρογγυλοποίηση δεν βοηθάει;

\ex \textbf{Τα πλακάκια του θείου.} Ο θείος Λευτέρης χρειάζεται 1.234 πλακάκια για τη βεράντα του. Πωλούνται σε κιβώτια των 100. Ο Στέλιος στρογγυλοποιεί στην πλησιέστερη εκατοντάδα και λέει «12 κιβώτια». Φτάνουν; Πόσα κιβώτια πρέπει να αγοράσει; Πόσα πλακάκια θα περισσέψουν;

\ex \textbf{Το βιβλίο της Μυρτώς.} Η Μυρτώ διαβάζει από τη σελίδα 47 ως τη σελίδα 112, και τις δύο μαζί. Πόσες σελίδες διαβάζει; Αν διαβάζει 6 σελίδες την ημέρα, φτάνουν 10 ημέρες;

\ex \textbf{Ο φράχτης.} Ο παππούς φράζει το χωράφι του κατά μήκος ενός δρόμου 20 μ. Βάζει έναν πάσσαλο κάθε 2 μ., και έναν στην αρχή και έναν στο τέλος. Πόσους πασσάλους θα βάλει;

\ex \textbf{Τα εισιτήρια.} Σε μια παράσταση τα εισιτήρια είναι αριθμημένα από το 120 ως το 399. Πόσα εισιτήρια υπάρχουν; Πόσα από αυτά έχουν περιττό αριθμό;

\ex \textbf{Ο άξονας του Ηλία.} Ο Ηλίας σχεδιάζει έναν άξονα που δείχνει το 0 και το 100 (δες το σχήμα). Ποιοι αριθμοί αντιστοιχούν στα σημεία Α, Β, Γ και Δ;

\begin{center}
\begin{tikzpicture}[x=1cm]
\draw[-{Stealth}] (-0.3,0)--(11.2,0);
\draw (0,-0.12)--(0,0.12) node[above,font=\small] {0};
\draw (2,-0.12)--(2,0.12) node[above,font=\small] {100};
\foreach \x/\n in {3/Α,5/Β,7.5/Γ,10/Δ}{\draw[gblue,thick] (\x,-0.15)--(\x,0.15); \node[below,font=\small\bfseries,gblue] at (\x,-0.15) {\n};}
\end{tikzpicture}
\end{center}

\zoni{Το σκέφτομαι}{μια ερώτηση που πάει πιο μακριά}

\ex Πόσες φορές γράφεται το ψηφίο 9, αν γράψεις όλους τους αριθμούς από το 1 ως το 100;

\ex Ποιος είναι ο μικρότερος και ποιος ο μεγαλύτερος φυσικός αριθμός που, στρογγυλοποιημένος στις χιλιάδες, δίνει 3.000;

\ex Ποιος είναι ο μεγαλύτερος και ποιος ο μικρότερος τετραψήφιος αριθμός που έχει όλα τα ψηφία του διαφορετικά;

\ex Ο Άρης λέει: «Θα βρω τον μεγαλύτερο φυσικό αριθμό και θα τον γράψω στον πίνακα». Η Ζωή λέει: «Δεν υπάρχει». Εξήγησε με λόγια γιατί έχει δίκιο η Ζωή.

\ex Ο Ηλίας δοκιμάζει στο κινητό: $7+5=12$, $3+9=12$. Ρωτάει τη Λίνα: «Αν προσθέσω δύο φυσικούς, βγαίνει πάντα φυσικός. Αν αφαιρέσω;» Η Λίνα γράφει $3-5$. Μπορεί να υπάρχει αυτό το αποτέλεσμα στους φυσικούς; Τι νομίζεις ότι θα χρειαστούμε;

\kefalaio{Πρόσθεση, αφαίρεση, πολλαπλασιασμός: οι ιδιότητες που κρύβονται}{Κεφάλαιο Α.1.2, «Πρόσθεση, αφαίρεση και πολλαπλασιασμός φυσικών αριθμών»}

\begin{anagke}
Η κυρία Αλεξάνδρα μπαίνει στην τάξη και λέει ότι κανένα κινητό δεν θα βγει από την τσάντα. Στον πίνακα γράφει: $98\times 7$. «Έχετε δέκα δευτερόλεπτα. Με το μυαλό».

Ο Στέλιος αρχίζει να γράφει στο χαρτί: οκτώ επί επτά, πενήντα έξι, κρατάω πέντε. Δεν προλαβαίνει. Ο Ηλίας σηκώνει το χέρι σε οκτώ δευτερόλεπτα: «Εξακόσια ογδόντα έξι». «Πώς;» ρωτάει η κυρία Αλεξάνδρα. «Το 98 είναι σχεδόν 100. Εκατό επί επτά, εφτακόσια. Έλειπαν δύο, δηλαδή δύο επί επτά, δεκατέσσερα. Εφτακόσια μείον δεκατέσσερα».

Η Ζωή κοιτάζει τον πίνακα. «Δηλαδή έσπασες το 98 σε 100 μείον 2, και πολλαπλασίασες το καθένα ξεχωριστά. Μπορείς να το κάνεις αυτό;» «Δεν ξέρω», λέει ο Ηλίας. «Βγαίνει σωστό». Η Λίνα συμπληρώνει: «Και γιατί να μη βγαίνει το ίδιο αν γράψω $7\times 98$;»

Δύο διαφορετικές ερωτήσεις, και οι δύο με την ίδια ρίζα. Τι ακριβώς επιτρέπεται να κάνουμε με τις πράξεις, και γιατί;
\end{anagke}

\begin{lexi}{αντιμεταθετική, προσεταιριστική, επιμεριστική}{αντιμεταθετική, προσεταιριστική, επιμεριστική}
Οι τρεις βαριές λέξεις έχουν απλό νόημα, αν τις σπάσεις.

\lexiword{Αντιμεταθετική}: αντι-μετατίθεμαι σημαίνει «αλλάζω θέση με τον άλλον». Οι δύο αριθμοί ανταλλάσσουν θέσεις και δεν αλλάζει τίποτα.

\lexiword{Προσεταιριστική}: από το «εταίρος», δηλαδή συνεργάτης, παρέα. Όταν προσθέτεις τρεις αριθμούς, διαλέγεις ποιοι δύο θα κάνουν πρώτα παρέα, με παρένθεση.

\lexiword{Επιμεριστική}: επι-μερίζω σημαίνει «μοιράζω σε μέρη». Ο αριθμός που πολλαπλασιάζει μοιράζεται σε καθένα από τα μέρη της παρένθεσης.
\end{lexi}

\begin{metafrasi}
Η σειρά των προσθετέων δεν αλλάζει το άθροισμα. & $\alpha+\beta=\beta+\alpha$\\
Όποιους δύο προσθέσω πρώτα, βγαίνει το ίδιο. & $(\alpha+\beta)+\gamma=\alpha+(\beta+\gamma)$\\
Αν προσθέσω 0, δεν αλλάζει τίποτα. & $\alpha+0=\alpha$\\
Είχα $Μ$, έδωσα $Α$ και μου έμειναν $\Delta$. & $Μ-Α=\Delta$, δηλαδή $Μ=Α+\Delta$\\
Το διπλάσιο του «τρία συν πέντε». & $2\cdot(3+5)=2\cdot 3+2\cdot 5$\\
Κάθε αριθμός επί 0. & $\alpha\cdot 0=0$\\
\end{metafrasi}

\tmima{Η θεωρία}

\tmimasmall{Οι τρεις πράξεις και τα ονόματά τους}

\begin{center}
\small\renewcommand{\arraystretch}{1.5}
\begin{tabular}{@{}l@{\hspace{1.2em}}l@{\hspace{1.2em}}l@{}}
\toprule
\textbf{Πράξη} & \textbf{Παράδειγμα} & \textbf{Ονόματα}\\
\midrule
Πρόσθεση & $13+5=18$ & 13 και 5: \emph{προσθετέοι}, 18: \emph{άθροισμα}\\
Αφαίρεση & $20-8=12$ & 20: \emph{μειωτέος}, 8: \emph{αφαιρετέος}, 12: \emph{διαφορά}\\
Πολλαπλασιασμός & $7\cdot 6=42$ & 7 και 6: \emph{παράγοντες}, 42: \emph{γινόμενο}\\
\bottomrule
\end{tabular}
\end{center}

\begin{idiotita}{πρόσθεση}
\begin{itemize}[leftmargin=1.2em,itemsep=0.2ex,topsep=2pt]
\item \textbf{Αντιμεταθετική}: $\alpha+\beta=\beta+\alpha$.
\item \textbf{Προσεταιριστική}: $(\alpha+\beta)+\gamma=\alpha+(\beta+\gamma)$. Γι' αυτό γράφουμε $\alpha+\beta+\gamma$ χωρίς παρένθεση.
\item Το 0 δεν αλλάζει το άθροισμα: $\alpha+0=0+\alpha=\alpha$.
\end{itemize}
\end{idiotita}

\begin{idiotita}{πολλαπλασιασμός}
\begin{itemize}[leftmargin=1.2em,itemsep=0.2ex,topsep=2pt]
\item \textbf{Αντιμεταθετική}: $\alpha\cdot\beta=\beta\cdot\alpha$. \quad \textbf{Προσεταιριστική}: $(\alpha\cdot\beta)\cdot\gamma=\alpha\cdot(\beta\cdot\gamma)$.
\item Το 1 δεν αλλάζει το γινόμενο: $\alpha\cdot 1=\alpha$. Το 0 το μηδενίζει: $\alpha\cdot 0=0$.
\item \textbf{Επιμεριστική} ως προς την πρόσθεση: $\alpha\cdot(\beta+\gamma)=\alpha\cdot\beta+\alpha\cdot\gamma$.
\item \textbf{Επιμεριστική} ως προς την αφαίρεση: $\alpha\cdot(\beta-\gamma)=\alpha\cdot\beta-\alpha\cdot\gamma$.
\end{itemize}
\end{idiotita}

\tmimasmall{Η αφαίρεση είναι ο αντίστροφος δρόμος}

Η αφαίρεση $Μ-Α=\Delta$ λέει: «ποιος αριθμός $\Delta$, αν τον προσθέσω στο $Α$, δίνει το $Μ$;». Γι' αυτό ισχύει πάντα $Μ=Α+\Delta$, και έτσι κάνουμε τη \textbf{δοκιμή}. Στους φυσικούς ο αφαιρετέος πρέπει να είναι μικρότερος ή ίσος από τον μειωτέο. Αλλιώς η πράξη δεν γίνεται. Η αφαίρεση δεν είναι ούτε αντιμεταθετική ούτε προσεταιριστική: $9-4\neq 4-9$ (το δεύτερο δεν υπάρχει στους φυσικούς) και $(10-5)-2=3$, ενώ $10-(5-2)=7$.

\tmimasmall{Η επιμεριστική με μια εικόνα}

Δύο ορθογώνια με το ίδιο ύψος $\gamma$ και πλάτη $\alpha$ και $\beta$ κολλάνε και φτιάχνουν ένα ορθογώνιο πλάτους $\alpha+\beta$. Το εμβαδό του είναι το άθροισμα των δύο: $(\alpha+\beta)\cdot\gamma=\alpha\gamma+\beta\gamma$.

\begin{center}
\begin{tikzpicture}[x=0.9cm,y=0.9cm]
\fill[gblueL] (0,0) rectangle (3,2); \fill[gwarm] (3,0) rectangle (5,2);
\draw[thick] (0,0) rectangle (5,2); \draw[thick] (3,0)--(3,2);
\node at (1.5,1) {$\alpha\cdot\gamma$}; \node at (4,1) {$\beta\cdot\gamma$};
\draw[<->] (0,-0.35)--(3,-0.35) node[midway,below] {$\alpha$};
\draw[<->] (3,-0.35)--(5,-0.35) node[midway,below] {$\beta$};
\draw[<->] (5.35,0)--(5.35,2) node[midway,right] {$\gamma$};
\end{tikzpicture}
\end{center}

\tmimasmall{Πολλαπλασιασμός επί 10, 100, 1.000, \dots}

\begin{kanonas}{επί 10, 100, 1.000}
Για να πολλαπλασιάσω έναν φυσικό επί 10, 100, 1.000, \dots, γράφω στο τέλος του τόσα μηδενικά όσα έχει ο 10, ο 100, ο 1.000, \dots Έτσι $421\cdot 100=42.100$ και $5\cdot 1.000=5.000$.
\end{kanonas}

\tmimasmall{Υπολογίζω με το μυαλό}

Οι ιδιότητες είναι εργαλεία για γρήγορους υπολογισμούς. Η αντιμεταθετική και η προσεταιριστική σου επιτρέπουν να ομαδεύεις αριθμούς που «κλείνουν» σε στρογγυλό: $4\cdot 17\cdot 25=17\cdot(4\cdot 25)=17\cdot 100=1.700$. Η επιμεριστική σου επιτρέπει να σπας έναν δύσκολο αριθμό σε εύκολα κομμάτια: $284\cdot 99=284\cdot(100-1)=28.400-284=28.116$.

\begin{prosoxi}
\begin{itemize}[leftmargin=1.2em,itemsep=0.3ex,topsep=2pt]
\item Η επιμεριστική πολλαπλασιάζει τον αριθμό με \textbf{κάθε} όρο της παρένθεσης. $5\cdot(2+3)=5\cdot 2+5\cdot 3=25$, όχι $5\cdot 2+3=13$.
\item Η αντιμεταθετική δεν ισχύει στην αφαίρεση: $9-4\neq 4-9$.
\item Αν στην αφαίρεση ένα ψηφίο του μειωτέου είναι μικρότερο από το αντίστοιχο του αφαιρετέου, δανείζεσαι μία δεκάδα από το διπλανό ψηφίο. Δεν γυρνάς τους αριθμούς ανάποδα.
\end{itemize}
\end{prosoxi}

\tmima{Λυμένα παραδείγματα}

\paradeigma{Το 98 επί 7, αναλυτικά}
\tixero Πρέπει να υπολογίσω το $98\cdot 7$ χωρίς κινητό.\par
\tithelo Το γινόμενο, με έναν τρόπο που να δικαιολογείται.\par
\possindeo Γράφω το 98 ως $100-2$ και χρησιμοποιώ την επιμεριστική ως προς την αφαίρεση: $98\cdot 7=(100-2)\cdot 7=100\cdot 7-2\cdot 7=700-14=686$.\par
\apantisi $686$.\par
\elegxos Το 98 είναι λίγο μικρότερο από το 100, άρα το γινόμενο πρέπει να είναι λίγο μικρότερο από το 700. Το 686 είναι.

\paradeigma{Η αντιμεταθετική σε δράση}
\tixero Πρέπει να υπολογίσω το $25\cdot 36$.\par
\tithelo Το γινόμενο.\par
\possindeo Το 36 ισούται με $4\cdot 9$. Άρα $25\cdot 36=25\cdot(4\cdot 9)=(25\cdot 4)\cdot 9=100\cdot 9=900$. Χρησιμοποίησα την προσεταιριστική, για να κάνω πρώτα το $25\cdot 4$ που κλείνει σε 100.\par
\apantisi $900$.\par
\elegxos Τέσσερα 25άρια είναι 100, άρα εννέα τετράδες από 25άρια είναι εννέα εκατοντάδες, δηλαδή 900. Σωστά.

\paradeigma{Πόσα είχε η Μυρτώ;}
\tixero Η Μυρτώ είχε κάποια χρήματα. Ξόδεψε 17 € στο κυλικείο και της έμειναν 38 €.\par
\tithelo Πόσα είχε στην αρχή.\par
\possindeo Μειωτέος είναι τα χρήματα που είχε, $Μ$. Αφαιρετέος είναι τα 17 και διαφορά τα 38. Από τη σχέση $Μ=Α+\Delta$ έχω $Μ=17+38=55$.\par
\apantisi 55 €.\par
\elegxos $55-17=38$. Σωστά.

\tmima{Ασκήσεις}

\zoni{Το κατάλαβα}{εφαρμογή του κανόνα}

\ex Συμπλήρωσε τα ονόματα. Στο $13+5=18$ το 13 και το 5 λέγονται \dots\ και το 18 \dots. Στο $20-8=12$ το 20 λέγεται \dots, το 8 \dots\ και το 12 \dots. Στο $7\cdot 6=42$ το 7 και το 6 λέγονται \dots\ και το 42 \dots.

\ex Ποια ιδιότητα χρησιμοποιείται σε κάθε ισότητα; α) $17+25=25+17$ \quad β) $(4+9)+1=4+(9+1)$ \quad γ) $6\cdot(10+3)=6\cdot 10+6\cdot 3$ \quad δ) $58\cdot 1=58$ \quad ε) $0\cdot 731=0$ \quad στ) $9+0=9$.

\ex Υπολόγισε: α) $52\cdot 100$ \quad β) $37\cdot 10$ \quad γ) $490\cdot 10.000$ \quad δ) $8\cdot 1.000$.

\ex Βρες τον αριθμό που λείπει: α) $17+\square=40$ \quad β) $\square+26=90$ \quad γ) $\square-15=32$ \quad δ) $80-\square=35$.

\ex Υπολόγισε με το μυαλό, ομαδεύοντας: α) $4\cdot 17\cdot 25$ \quad β) $8\cdot 39\cdot 125$ \quad γ) $12+47+88+53$ \quad δ) $5\cdot 73\cdot 2$.

\ex Υπολόγισε με την επιμεριστική ιδιότητα: α) $89\cdot 7+89\cdot 3$ \quad β) $23\cdot 49+77\cdot 49$ \quad γ) $76\cdot 13-76\cdot 3$ \quad δ) $284\cdot 99$ \quad ε) $35\cdot 102$.

\ex Σωστό ή λάθος; α) $9-4=4-9$ \quad β) $(20-8)-2=20-(8-2)$ \quad γ) $5\cdot(2+3)=5\cdot 2+3$ \quad δ) $7\cdot 0=7$ \quad ε) $12-12=0$ \quad στ) $\alpha\cdot 1=\alpha$ για κάθε φυσικό $\alpha$.

\ex Βρες τα αποτελέσματα: α) $157+33$ \quad β) $122+25+78$ \quad γ) $785-323$ \quad δ) $7.321-4.595$ \quad ε) $60-(18-2)$.

\ex Ταίριαξε κάθε παράσταση με την τιμή της: $1+2+3+4$, \ $1+2+34$, \ $12+34$, \ $1\cdot 2\cdot 3\cdot 4$. Τιμές: 10, 24, 37, 46.

\zoni{Το χρησιμοποιώ}{προβλήματα από τη ζωή}

\ex \textbf{Το κυλικείο.} Το σχολικό κυλικείο πουλάει 98 τοστ κάθε ημέρα. Πόσα τοστ πουλάει σε 5 ημέρες; Υπολόγισε το $98\cdot 5$ με το μυαλό, γράφοντας το 98 ως $100-2$.

\ex \textbf{Οι πινακίδες.} Στον δρόμο για το χωριό του παππού οι πινακίδες γράφουν χιλιόμετρα από την αφετηρία: γέφυρα 47, βενζινάδικο 112, χωριό 175, θάλασσα 230. Πόσα χιλιόμετρα είναι από την αφετηρία ως τη γέφυρα, και μετά από κάθε πινακίδα ως την επόμενη; Πόσο απέχει το βενζινάδικο από τη θάλασσα;

\ex \textbf{Τα αυτοκόλλητα.} Η Ζωή έδωσε 17 αυτοκόλλητα στη Λίνα και της έμειναν 38. Πόσα είχε στην αρχή; Έλεγξε την απάντησή σου με μια αφαίρεση.

\ex \textbf{Η αφαίρεση του Άρη.} Ο Άρης κάνει την αφαίρεση $82-47$ και βγάζει 45. Ο τρόπος του: σε κάθε στήλη αφαιρεί το μικρότερο ψηφίο από το μεγαλύτερο. Γιατί είναι λάθος; Πώς θα το έλεγχες; Βρες το σωστό αποτέλεσμα.

\ex \textbf{Το πάτωμα του θείου.} Ο θείος Λευτέρης αγοράζει 48 κιβώτια πλακάκια με 25 πλακάκια το καθένα. Πόσα πλακάκια είναι; Υπολόγισε το $48\cdot 25$ γράφοντας το 48 ως $12\cdot 4$.

\ex \textbf{Ο κήπος της γιαγιάς.} Ο ορθογώνιος κήπος της γιαγιάς Ρήνης έχει μήκος 13 μ. και πλάτος 8 μ. Σπάσε το 13 σε $10+3$ και υπολόγισε το εμβαδό του σαν άθροισμα δύο ορθογωνίων. Φτιάξε ένα σχήμα.

\ex \textbf{Τα παπούτσια.} Σε ένα κατάστημα τα αθλητικά παπούτσια κοστίζουν 99 € το ζευγάρι. Ο Άρης και οι δύο συμπαίκτες του παίρνουν από ένα. Πόσο πληρώνουν συνολικά; Υπολόγισέ το με το μυαλό.

\zoni{Το σκέφτομαι}{μια ερώτηση που πάει πιο μακριά}

\ex \textbf{Ο Γκάους.} Όταν ο Γκάους ήταν παιδί, ο δάσκαλός του ζήτησε να προσθέσουν τους αριθμούς από το 1 ως το 100. Εκείνος τους έγραψε σε ζευγάρια: $(1+100)+(2+99)+\dots$ και βρήκε την απάντηση σε δευτερόλεπτα. Βρες την κι εσύ. Πόσα ζευγάρια υπάρχουν και πόσο κάνει το καθένα; Με την ίδια ιδέα βρες το $1+2+3+\dots+10$ και έλεγξε το αποτέλεσμα με πρόσθεση.

\ex Δύο αριθμοί έχουν άθροισμα 12 και διαφορά 2. Ποιοι είναι;

\ex Η Μυρτώ λέει ότι $4\cdot 0=0$. Ο Στέλιος ρωτάει «Γιατί;». Εξήγησε με την εικόνα «4 σακούλες με 0 καραμέλες η καθεμία».

\ex Η Λίνα ρώτησε αν το $3\cdot 5$ και το $5\cdot 3$ είναι το ίδιο. Σχεδίασε ένα ορθογώνιο με 3 σειρές των 5 τετραγώνων και εξήγησε γιατί, αν το γυρίσεις στο πλάι, βλέπεις 5 σειρές των 3. Είναι «το ίδιο» ή απλώς «ίδιο αποτέλεσμα»;

\ex Ο Ηλίας δοκιμάζει στο τετράδιό του τα $(10-5)-2$ και $10-(5-2)$. Βρες και τις δύο τιμές. Τι δείχνει αυτό για την αφαίρεση; Γι' αυτό στο $10-5-2$ συμφωνούμε να εκτελούμε τις πράξεις από αριστερά προς τα δεξιά.

\ex Πόσο κάνει το $3\cdot 3\cdot 3\cdot 3\cdot 3\cdot 3$; Πόσα ψηφία έχει το αποτέλεσμα; Σκέψου πώς θα το έγραφες πιο σύντομα. Αν το καταφέρεις, θα έχεις βρει το θέμα του επόμενου κεφαλαίου.

\kefalaio{Δυνάμεις και προτεραιότητα: ποιος πάει πρώτος}{Κεφάλαιο Α.1.3, «Δυνάμεις φυσικών αριθμών»}

\begin{anagke}
Παρασκευή απόγευμα. Η Λίνα φτιάχνει κωδικό για το καινούργιο της κινητό: τέσσερα ψηφία. «Πόσοι διαφορετικοί κωδικοί υπάρχουν;» ρωτάει η Ζωή. Η Λίνα σκέφτεται δυνατά: «Για το πρώτο ψηφίο έχω 10 επιλογές, από το 0 ως το 9. Για το δεύτερο πάλι 10, και για το τρίτο, και για το τέταρτο». Ο Ηλίας γράφει $10\cdot 10\cdot 10\cdot 10$ και βρίσκει 10.000. «Και αν ο κωδικός είχε είκοσι ψηφία;» ρωτάει. Γράφει είκοσι δεκάρια στη σειρά και χάνει το μέτρημα στη μέση.

Ο Άρης, που βιάζεται πάντα, γράφει σε ένα χαρτάκι $2+3\cdot 4$ και ρωτάει τη Ζωή πόσο κάνει. «Δεκατέσσερα», λέει εκείνη. «Είκοσι!» φωνάζει ο Άρης. «Δύο και τρία πέντε, επί τέσσερα, είκοσι». Ο Θοδωρής, που περνάει από δίπλα, τους κοιτάει και λέει μόνο: «Και οι δύο έχετε ένα κομπιουτεράκι στην τσέπη. Δείτε τι βγάζει». Το κομπιουτεράκι δίνει 14.

Δύο προβλήματα. Το πρώτο: πώς γράφεται σύντομα ο ίδιος παράγοντας πολλές φορές; Το δεύτερο: όταν μια παράσταση έχει πολλές πράξεις, ποια γίνεται πρώτη, και ποιος το αποφασίζει;
\end{anagke}

\begin{lexi}{δύναμη, εκθέτης, βάση}{δύναμη, εκθέτης}
Η λέξη \lexiword{δύναμη} στα μαθηματικά δεν έχει σχέση με τα μπράτσα. Ο αριθμός «έχει τη δύναμη» να πολλαπλασιάζεται με τον εαυτό του πολλές φορές. Το $2^{10}$ διαβάζεται «το δύο στη δέκατη».

Ο \lexiword{εκθέτης} είναι το μικρό νούμερο πάνω δεξιά. Το «εκθέτω» σημαίνει «δείχνω»: ο εκθέτης δείχνει πόσες φορές μπαίνει ο παράγοντας. Ο αριθμός που πολλαπλασιάζεται λέγεται \lexiword{βάση}, γιατί πάνω του «στηρίζεται» το σύμβολο.
\end{lexi}

\begin{metafrasi}
Το 2 πολλαπλασιασμένο με τον εαυτό του 10 φορές. & $2^{10}=2\cdot 2\cdot\dots\cdot 2$ (10 παράγοντες)\\
Το τετράγωνο του 5. & $5^2=5\cdot 5=25$\\
Ο κύβος του 3. & $3^3=3\cdot 3\cdot 3=27$\\
Το 7 στην πρώτη. & $7^1=7$\\
Ένα εκατομμύριο, με δύναμη του 10. & $10^6=1.000.000$\\
Πρώτα οι δυνάμεις, μετά οι πολλαπλασιασμοί. & $2+3\cdot 4^2=2+3\cdot 16=50$\\
\end{metafrasi}

\tmima{Η θεωρία}

\begin{orismos}{δύναμη}
Το γινόμενο $\alpha\cdot\alpha\cdot\alpha\cdots\alpha$ με $\nu$ παράγοντες ίσους με το $\alpha$ λέγεται \textbf{δύναμη} του $\alpha$ στη $\nu$ και γράφεται $\alpha^{\nu}$. Ο $\alpha$ είναι η \textbf{βάση} και ο $\nu$ ο \textbf{εκθέτης}.
\end{orismos}

Η δεύτερη δύναμη $\alpha^2$ λέγεται και \textbf{τετράγωνο} του $\alpha$ (είναι το εμβαδό τετραγώνου με πλευρά $\alpha$). Η τρίτη δύναμη $\alpha^3$ λέγεται και \textbf{κύβος} του $\alpha$ (είναι ο όγκος κύβου με ακμή $\alpha$). Η πρώτη δύναμη $\alpha^1$ είναι ο ίδιος ο αριθμός $\alpha$. Οι δυνάμεις του 1 είναι όλες 1: $1^{\nu}=1$.

\begin{prosoxi}
Το $2^3$ δεν είναι $2\cdot 3=6$ αλλά $2\cdot 2\cdot 2=8$. Ο εκθέτης δείχνει πόσες φορές γράφεται η βάση, όχι επί πόσο την πολλαπλασιάζουμε.
\end{prosoxi}

\tmimasmall{Οι δυνάμεις του 10}

Οι δυνάμεις του 10 έχουν τόσα μηδενικά όσος ο εκθέτης:
$$10^2=100,\qquad 10^3=1.000,\qquad 10^6=1.000.000.$$
Γι' αυτό η δεκαδική μας γραφή αναλύεται εύκολα. Για παράδειγμα:
$$7.604=7\cdot 10^3+6\cdot 10^2+0\cdot 10+4.$$

\tmimasmall{Αριθμητική παράσταση και προτεραιότητα}

Μια \textbf{αριθμητική παράσταση} είναι μια σειρά αριθμών που συνδέονται με σύμβολα πράξεων. Ο αριθμός που βγαίνει όταν κάνουμε όλες τις πράξεις λέγεται \textbf{τιμή} της. Για να βγάζουμε όλοι την ίδια τιμή, όλος ο κόσμος έχει συμφωνήσει σε μια σειρά.

\begin{kanonas}{προτεραιότητα των πράξεων}
\begin{enumerate}[leftmargin=1.6em,itemsep=0.2ex,topsep=2pt]
\item Πρώτα οι πράξεις μέσα στις \textbf{παρενθέσεις}, με την ίδια σειρά που θα βρεις παρακάτω.
\item Μετά οι \textbf{δυνάμεις}.
\item Μετά οι \textbf{πολλαπλασιασμοί και οι διαιρέσεις}, από αριστερά προς τα δεξιά.
\item Τελευταίες οι \textbf{προσθέσεις και οι αφαιρέσεις}, από αριστερά προς τα δεξιά.
\end{enumerate}
\end{kanonas}

Η συμφωνία αυτή δεν είναι «φυσικός νόμος». Είναι συμφωνία, όπως ο κανόνας ότι στον δρόμο πάμε δεξιά. Αν δεν υπήρχε, το $2+3\cdot 4$ θα είχε δύο απαντήσεις (14 και 20) και δύο μηχανικοί δεν θα συμφωνούσαν ποτέ σε έναν υπολογισμό. Με τη συμφωνία η απάντηση είναι μία, 14, και η Ζωή έχει δίκιο. Αν θέλουμε το 20, γράφουμε $(2+3)\cdot 4$: η παρένθεση λέει ρητά «αυτό πρώτα».

\begin{prosoxi}
\begin{itemize}[leftmargin=1.2em,itemsep=0.3ex,topsep=2pt]
\item Στο $2\cdot 3^2$ ο εκθέτης αφορά μόνο το 3: $2\cdot 9=18$. Για να τετραγωνίσεις όλο το γινόμενο, γράφεις $(2\cdot 3)^2=36$.
\item Οι πράξεις του ίδιου επιπέδου γίνονται από αριστερά προς τα δεξιά: $24:4:2=6:2=3$, όχι $24:(4:2)=12$.
\end{itemize}
\end{prosoxi}

\tmima{Λυμένα παραδείγματα}

\paradeigma{Οι κωδικοί της Λίνας}
\tixero Ο κωδικός έχει 4 ψηφία και κάθε ψηφίο έχει 10 επιλογές (0 ως 9).\par
\tithelo Πόσοι είναι οι δυνατοί κωδικοί, και πώς γράφεται σύντομα ο υπολογισμός;\par
\possindeo Για κάθε ψηφίο έχω 10 επιλογές, και οι επιλογές συνδυάζονται, άρα πολλαπλασιάζω: $10\cdot 10\cdot 10\cdot 10$. Είναι τέσσερα δεκάρια, δηλαδή $10^4$. Υπολογίζω βήμα βήμα: $10\cdot 10=100$, $100\cdot 10=1.000$, $1.000\cdot 10=10.000$.\par
\apantisi $10^4=10.000$ κωδικοί.\par
\elegxos Οι κωδικοί είναι οι αριθμοί από το 0000 ως το 9999. Από το κεφάλαιο 1: $9.999-0+1=10.000$. Βγήκε το ίδιο.

\paradeigma{Η παράσταση με πολλές πράξεις}
\tixero Η παράσταση $(2\cdot 5)^4+4\cdot(3+2)^2$.\par
\tithelo Την τιμή της.\par
\possindeo Πρώτα οι παρενθέσεις: $2\cdot 5=10$ και $3+2=5$. Έχω $10^4+4\cdot 5^2$. Μετά οι δυνάμεις: $10^4=10.000$ και $5^2=25$. Μετά ο πολλαπλασιασμός: $4\cdot 25=100$. Τέλος η πρόσθεση: $10.000+100=10.100$.\par
\apantisi $10.100$.\par
\elegxos Το $10^4$ έχει 4 μηδενικά, 10.000. Το άθροισμα με το 100 δίνει 10.100, όχι κάτι πιο μικρό από το 10.000. Σωστά.

\paradeigma{Ο αριθμός του Ηλία}
\tixero Ο Ηλίας γράφει ένα αριθμό σε ανάπτυγμα: $5\cdot 10^4+3\cdot 10^2+8$.\par
\tithelo Ποιος αριθμός είναι;\par
\possindeo $5\cdot 10^4=5\cdot 10.000=50.000$. $3\cdot 10^2=300$. Άρα $50.000+300+8=50.308$.\par
\apantisi $50.308$.\par
\elegxos Λείπουν οι χιλιάδες και οι δεκάδες, που σημαίνουν 0 στη θέση τους: 50.308. Τα μηδενικά έχουν μπει σωστά.

\tmima{Ασκήσεις}

\zoni{Το κατάλαβα}{εφαρμογή του κανόνα}

\ex Γράψε ως δύναμη: α) $5\cdot 5\cdot 5$ \quad β) $2\cdot 2\cdot 2\cdot 2\cdot 2$ \quad γ) $10\cdot 10\cdot 10\cdot 10$ \quad δ) $7\cdot 7$ \quad ε) $9$.

\ex Γράψε ως γινόμενο και υπολόγισε: α) $2^4$ \quad β) $3^3$ \quad γ) $5^2$ \quad δ) $10^5$ \quad ε) $1^{100}$ \quad στ) $6^1$.

\ex Ποια είναι η βάση και ποιος ο εκθέτης στα $4^5$, $12^2$, $1^{9}$ και $10^8$;

\ex Υπολόγισε το τετράγωνο και τον κύβο των αριθμών 2, 3, 4, 5 και 10. Φτιάξε πίνακα.

\ex Υπολόγισε την τιμή: α) $12+3\cdot 4$ \quad β) $(12+3)\cdot 4$ \quad γ) $24:4:2$ \quad δ) $100-5\cdot 4$ \quad ε) $2+3^2$ \quad στ) $(2+3)^2$.

\ex Βρες την τιμή: α) $(2\cdot 5)^4+4\cdot(3+2)^2$ \quad β) $(2+3)^3-8\cdot 3^2$ \quad γ) $3\cdot(4+2)^2-10\cdot 3$.

\ex Γράψε το ανάπτυγμα με δυνάμεις του 10: α) 7.604 \quad β) 53.280 \quad γ) 400.065.

\ex Γράψε με ψηφία τον αριθμό: α) $3\cdot 10^3+2\cdot 10^2+5$ \quad β) $9\cdot 10^5+4\cdot 10$ \quad γ) $10^6+10^2$.

\ex Σωστό ή λάθος; α) $2^3=6$ \quad β) $10^4=10.000$ \quad γ) $5^2=2\cdot 5$ \quad δ) $1^{57}=1$ \quad ε) $3+2\cdot 5=25$ \quad στ) $(1+1)^2=1^2+1^2$.

\zoni{Το χρησιμοποιώ}{προβλήματα από τη ζωή}

\ex \textbf{Το λουκέτο.} Ένα λουκέτο έχει τροχούς με ψηφία από το 0 ως το 9. Πόσους διαφορετικούς κωδικούς έχει αν έχει 3 τροχούς; Αν έχει 4; Αν έχει 5; Γράψε κάθε απάντηση και ως δύναμη του 10.

\ex \textbf{Οι διακόπτες.} Στην αποθήκη του θείου Λευτέρη ο πίνακας έχει διακόπτες που ο καθένας έχει δύο θέσεις: ανοιχτός ή κλειστός. Πόσους διαφορετικούς συνδυασμούς θέσεων έχουν 3 διακόπτες; Πόσους 5; Πόσους 10; Γράψε κάθε απάντηση και ως δύναμη του 2.

\ex \textbf{Ο κύβος του θείου.} Ο θείος Λευτέρης έχει δοχείο σε σχήμα κύβου με ακμή 4 δεκατόμετρα. Πόσα κυβικά δεκατόμετρα χωράει; (Θυμήσου ότι ο όγκος κύβου είναι $\alpha^3$.)

\ex \textbf{Ο κήπος.} Ο τετράγωνος κήπος της γιαγιάς Ρήνης έχει πλευρά 12 μ. Βρες το εμβαδό του. Αν διπλασιαστεί η πλευρά, ο κήπος γίνεται διπλάσιος σε εμβαδό; Υπολόγισε.

\ex \textbf{Τα σκακιστικά τετράγωνα.} Η σκακιέρα έχει 8 γραμμές και 8 στήλες. Πόσα τετράγωνα έχει; Γράψε τον υπολογισμό με δύναμη.

\ex \textbf{Ο λογαριασμός του Άρη.} Ο Άρης αγοράζει 3 μπάλες των 8 € και ένα φανελάκι των 12 €. Γράφει $3\cdot 8+12$ και βρίσκει 36. Η Μυρτώ γράφει $3\cdot(8+12)$. Ποια παράσταση ταιριάζει στο πρόβλημα και πόσο πληρώνει τελικά ο Άρης;

\ex \textbf{Τρεις απαντήσεις.} Τρεις φίλοι υπολογίζουν την παράσταση $4\cdot(7+7\cdot 9)+20$. Ο Άρης βρίσκει 335, η Ζωή 300 και ο Ηλίας 524. Ποιος έχει δίκιο; Βρες πού έκανε λάθος καθένας από τους άλλους δύο.

\zoni{Το σκέφτομαι}{μια ερώτηση που πάει πιο μακριά}

\ex Πόσο κάνει το $2^3\cdot 2^2$; Γράψε το ως γινόμενο παραγόντων (όλα τα δυάρια) και μέτρα. Ανακαλύπτεις έναν γρήγορο κανόνα για τον πολλαπλασιασμό δυνάμεων με την ίδια βάση;

\ex Συγκρίνε τα $2^5$ και $5^2$, και τα $2^4$ και $4^2$. Μετά βρες έναν αριθμό $\alpha$ διαφορετικό από το 2 για τον οποίο το $\alpha^2$ είναι ίσο με το $2\alpha$.

\ex Τοποθέτησε παρενθέσεις στην παράσταση $2+3\cdot 4+5$ με όσους τρόπους μπορείς και βρες όλες τις διαφορετικές τιμές που μπορεί να πάρει.

\ex Βάλε τα σύμβολα $+,-,\cdot$ (ή παρενθέσεις) ανάμεσα στα ψηφία $1\ \ 2\ \ 3\ \ 4$ ώστε το αποτέλεσμα να είναι 10. Βρες τουλάχιστον δύο τρόπους.

\ex Ο Ηλίας λέει: «Το $10^0$ πρέπει να είναι 0, γιατί δεν πολλαπλασιάζεις το 10 καμία φορά». Η Ζωή λέει: «Μάλλον είναι 1». Δες το μοτίβο $10^3=1.000$, $10^2=100$, $10^1=10$ και σκέψου τι πρέπει να συνεχίσει. Ποιος νομίζεις ότι έχει δίκιο; (Θα το μάθεις στο σχολείο σε λίγους μήνες.)

\kefalaio{Ευκλείδεια διαίρεση: το υπόλοιπο που μένει}{Κεφάλαιο Α.1.4, «Ευκλείδεια διαίρεση – Διαιρετότητα»}

\begin{anagke}
Το σχολείο οργανώνει την πρώτη του εκδρομή. Θα πάνε και τα 387 παιδιά, με λεωφορεία των 49 θέσεων. Η κυρία Αλεξάνδρα ζητάει από την τάξη να βρει πόσα λεωφορεία χρειάζονται.

Ο Στέλιος, που αγαπάει το κομπιουτεράκι, κάνει $387:49$ και η οθόνη γράφει $7{,}897959\dots$. «Εφτά κόμμα εννιά!» λέει περήφανος. «Άρα εφτά λεωφορεία». Η Μυρτώ τον κοιτάει. «Και τα υπόλοιπα παιδιά; Μένουν στην αυλή;» Ο Στέλιος ξαναδιαβάζει την οθόνη. «Μα το 0,897 τι είναι; Δεν υπάρχει 0,897 λεωφορείου». Η Ζωή κάνει τη διαίρεση στο χαρτί, όπως τη μάθαμε: πηλίκο 7 και κάτι περισσεύει. «Αυτό που περισσεύει είναι παιδιά, όχι κόμματα. Εκεί βρίσκεται η απάντηση».

Σε μια διαίρεση, πολλές φορές, δεν ενδιαφέρει μόνο πόσες φορές χωράει ο ένας αριθμός στον άλλον. Ενδιαφέρει και τι περισσεύει. Πώς γράφεται αυτό με κανόνα;
\end{anagke}

\begin{lexi}{πηλίκο, υπόλοιπο}{πηλίκο, υπόλοιπο}
Το \lexiword{πηλίκο} έρχεται από το αρχαίο «πηλίκος», που σημαίνει «πόσο μεγάλος». Η διαίρεση απαντάει στο «πόσο μεγάλο μερίδιο παίρνει καθένας». Το \lexiword{υπόλοιπο} είναι κυριολεκτικά αυτό που μένει: υπό-λοιπο, το «λοιπό» που δεν μοιράστηκε.

Η \lexiword{Ευκλείδεια διαίρεση} πήρε το όνομά της από τον Ευκλείδη, τον Έλληνα μαθηματικό που έζησε στην Αλεξάνδρεια περίπου το 300 π.Χ. και έγραψε τα «Στοιχεία», το βιβλίο όπου συγκέντρωσε όλα τα μαθηματικά της εποχής του. Η Ευκλείδεια διαίρεση είναι η διαίρεση με πηλίκο και υπόλοιπο, όπως τη μαθαίνουμε στο δημοτικό.
\end{lexi}

\begin{metafrasi}
Μοιράζω 43 μπάλες σε ομάδες των 7: γίνονται 6 ομάδες και μένει 1 μπάλα. & $43=7\cdot 6+1$\\
Δεν περισσεύει τίποτα (τέλεια διαίρεση). & $\Delta=\delta\cdot\pi$, υπόλοιπο 0\\
Το υπόλοιπο είναι πάντα μικρότερο από τον διαιρέτη. & $\upsilon<\delta$\\
Στη διαίρεση με το 8 τα πιθανά υπόλοιπα. & $0,\,1,2,\,3,4,\,5,6,\,7$\\
Δεν διαιρούμε ποτέ με το μηδέν. & $\delta\neq 0$\\
\end{metafrasi}

\tmima{Η θεωρία}

\begin{orismos}{Ευκλείδεια διαίρεση}
Όταν δίνονται δύο φυσικοί $\Delta$ και $\delta$ (με $\delta\neq 0$), υπάρχουν δύο φυσικοί $\pi$ και $\upsilon$ με
$$\Delta=\delta\cdot\pi+\upsilon,\qquad \upsilon<\delta.$$
Ο $\Delta$ λέγεται \textbf{διαιρετέος}, ο $\delta$ \textbf{διαιρέτης}, ο $\pi$ \textbf{πηλίκο} και ο $\upsilon$ \textbf{υπόλοιπο}.
\end{orismos}

Η συνθήκη $\upsilon<\delta$ είναι το κλειδί. Αν το υπόλοιπο ήταν ίσο ή μεγαλύτερο από τον διαιρέτη, θα χωρούσε άλλη μία φορά ο διαιρέτης και θα μπορούσαμε να αυξήσουμε το πηλίκο. Η \textbf{δοκιμή} της διαίρεσης είναι η ίδια η ισότητα: $\delta\cdot\pi+\upsilon$ πρέπει να δίνει τον $\Delta$.

Όταν $\upsilon=0$ έχουμε \textbf{τέλεια διαίρεση}: $\Delta=\delta\cdot\pi$. Τότε λέμε ότι ο $\Delta$ \textbf{διαιρείται} με τον $\delta$, και ότι ο $\delta$ \textbf{διαιρεί} τον $\Delta$. Στην τέλεια διαίρεση η διαίρεση είναι η αντίστροφη πράξη του πολλαπλασιασμού: αν $\Delta=\delta\cdot\pi$, τότε $\Delta:\delta=\pi$ και $\Delta:\pi=\delta$.

\begin{idiotita}{διαίρεση}
\begin{itemize}[leftmargin=1.2em,itemsep=0.2ex,topsep=2pt]
\item Ο διαιρέτης δεν μπορεί να είναι 0: δεν υπάρχει $\Delta:0$.
\item $\alpha:\alpha=1$ (αν $\alpha\neq 0$). \quad $\alpha:1=\alpha$. \quad $0:\alpha=0$ (αν $\alpha\neq 0$).
\end{itemize}
\end{idiotita}

Γιατί δεν διαιρούμε με το 0; Αν ήταν $6:0=x$, τότε θα έπρεπε $6=0\cdot x$. Μα κάθε αριθμός επί 0 κάνει 0, ποτέ 6. Δεν υπάρχει κανένας τέτοιος $x$.

\begin{prosoxi}
\begin{itemize}[leftmargin=1.2em,itemsep=0.3ex,topsep=2pt]
\item Η ισότητα $\Delta=\delta\cdot\pi+\upsilon$ δεν αρκεί για να είναι Ευκλείδεια διαίρεση. Πρέπει και $\upsilon<\delta$. Το $1.345=59\cdot 21+106$ είναι σωστή ισότητα, αλλά το 106 είναι μεγαλύτερο από το 59 και από το 21: δεν είναι διαίρεση.
\item Το πηλίκο δεν είναι πάντα η απάντηση σε ένα πρόβλημα. Για 387 παιδιά και λεωφορεία των 49 θέσεων, το πηλίκο είναι 7, αλλά χρειάζονται 8 λεωφορεία, γιατί το υπόλοιπο παιδιά θέλει θέσεις. Διάβασε τι ρωτάει το πρόβλημα.
\end{itemize}
\end{prosoxi}

\tmima{Λυμένα παραδείγματα}

\paradeigma{Τα λεωφορεία}
\tixero 387 παιδιά, λεωφορεία με 49 θέσεις.\par
\tithelo Πόσα λεωφορεία χρειάζονται και πόσες θέσεις θα μείνουν άδειες.\par
\possindeo Κάνω την Ευκλείδεια διαίρεση $387:49$. Το $49\cdot 7=343$ και το $49\cdot 8=392$ είναι πάνω από το 387, άρα πηλίκο 7 και υπόλοιπο $387-343=44$. Γράφω $387=49\cdot 7+44$. Επτά λεωφορεία γεμίζουν, και 44 παιδιά περισσεύουν. Χρειάζεται ένα όγδοο λεωφορείο γι' αυτά. Σε αυτό θα μείνουν $49-44=5$ θέσεις άδειες.\par
\apantisi 8 λεωφορεία, και 5 άδειες θέσεις.\par
\elegxos Η δοκιμή: $49\cdot 7+44=343+44=387$. Και $44<49$, όπως πρέπει.

\paradeigma{Ποιες ισότητες είναι διαιρέσεις;}
\tixero Οι ισότητες α) $120=28\cdot 4+8$, β) $1.345=59\cdot 21+106$, γ) $374=8\cdot 46+6$.\par
\tithelo Σε ποιες περιπτώσεις είναι Ευκλείδεια διαίρεση και με ποιον διαιρέτη.\par
\possindeo Ελέγχω ότι η ισότητα ισχύει και ότι το υπόλοιπο είναι μικρότερο από τον διαιρέτη.\par
α) Το υπόλοιπο 8 είναι μικρότερο από το 28, όχι από το 4. Είναι διαίρεση με διαιρέτη το 28 και πηλίκο 4.\par
β) Το υπόλοιπο 106 είναι μεγαλύτερο και από το 59 και από το 21. Δεν είναι διαίρεση.\par
γ) Το υπόλοιπο 6 είναι μικρότερο και από το 8 και από το 46. Είναι διαίρεση και με διαιρέτη το 8 (πηλίκο 46) και με διαιρέτη το 46 (πηλίκο 8).\par
\apantisi Μόνο οι (α) και (γ).\par
\elegxos Για το (β): $59\cdot 21=1.239$ και $1.345-1.239=106$. Ναι, η ισότητα ισχύει, αλλά το 106 περισσεύει πολύ: ο 59 χωράει ακόμα μία φορά στο 106. Άρα το πηλίκο δεν ήταν 21.

\paradeigma{Η μέρα της εβδομάδας}
\tixero Σήμερα είναι Τρίτη.\par
\tithelo Τι μέρα θα είναι μετά από 247 ημέρες.\par
\possindeo Η εβδομάδα έχει 7 μέρες, και κάθε 7 μέρες γυρίζουμε στην ίδια μέρα. Διαιρώ: $247=7\cdot 35+2$. Οι 35 ολόκληρες εβδομάδες δεν αλλάζουν τη μέρα. Μένουν 2 μέρες: Τετάρτη, Πέμπτη.\par
\apantisi Πέμπτη.\par
\elegxos $7\cdot 35=245$, και μετά από 245 ημέρες είναι πάλι Τρίτη. Δύο μέρες μετά είναι Πέμπτη. Το υπόλοιπο μας είπε πόσο προχωράμε μέσα στην εβδομάδα.

\tmima{Ασκήσεις}

\zoni{Το κατάλαβα}{εφαρμογή του κανόνα}

\ex Κάνε τις διαιρέσεις και τη δοκιμή τους: α) $4.002:69$ \quad β) $1.445:17$ \quad γ) $925:37$ \quad δ) $3.621:213$ \quad ε) $35.280:2.940$ \quad στ) $5.082:77$.

\ex Βρες πηλίκο και υπόλοιπο και γράψε την ισότητα $\Delta=\delta\cdot\pi+\upsilon$: α) $43:7$ \quad β) $100:9$ \quad γ) $245:12$ \quad δ) $1.000:33$ \quad ε) $59:60$ \quad στ) $0:5$.

\ex Ποιες από τις παρακάτω ισότητες εκφράζουν Ευκλείδεια διαίρεση; Με ποιον διαιρέτη; α) $125=35\cdot 3+20$ \quad β) $762=38\cdot 19+40$ \quad γ) $1.500=42\cdot 35+30$ \quad δ) $300=18\cdot 16+12$.

\ex Αν $\nu$ είναι ένας φυσικός αριθμός, ποια υπόλοιπα μπορεί να έχει η διαίρεση $\nu:8$;

\ex Ένας αριθμός διαιρούμενος με το 9 δίνει πηλίκο 73 και υπόλοιπο 4. Ποιος είναι ο αριθμός;

\ex Συμπλήρωσε: α) $\delta=12$, $\pi=7$, $\upsilon=5$: βρες τον $\Delta$. \quad β) $\Delta=100$, $\delta=8$: βρες $\pi$ και $\upsilon$. \quad γ) $\Delta=59$, $\pi=4$, $\upsilon=3$: βρες τον $\delta$.

\ex Σωστό ή λάθος; α) Στη διαίρεση με το 6 το υπόλοιπο μπορεί να είναι 6. \quad β) $0:7=0$. \quad γ) $7:0=0$. \quad δ) $\alpha:\alpha=1$ για κάθε $\alpha\neq 0$. \quad ε) $17:1=1$. \quad στ) Όταν το υπόλοιπο είναι 0, η διαίρεση λέγεται τέλεια.

\ex Ο Στέλιος γράφει $85=12\cdot 6+13$ και λέει «είναι διαίρεση, γιατί η ισότητα ισχύει». Ισχύει η ισότητα; Γιατί παρ' όλα αυτά δεν είναι Ευκλείδεια διαίρεση; Διόρθωσέ την.

\zoni{Το χρησιμοποιώ}{προβλήματα από τη ζωή}

\ex \textbf{Η εκδρομή.} Στην εκδρομή του σχολείου πάνε 387 παιδιά και κάθε λεωφορείο έχει 49 θέσεις. Πόσα λεωφορεία χρειάζονται; Πόσες θέσεις μένουν κενές στο τελευταίο; Αν η εταιρεία έχει μόνο 7 λεωφορεία, πόσα παιδιά θα μείνουν στο σχολείο;

\ex \textbf{Οι μέρες.} Σήμερα είναι Τρίτη. Τι μέρα θα είναι μετά από 150 ημέρες; Τι μέρα ήταν πριν από 8 ημέρες;

\ex \textbf{Οι γλάστρες του παππού.} Ο παππούς έχει 1.000 σπόρους και φυτεύει 33 σε κάθε γλάστρα. Πόσες γλάστρες γεμίζει; Πόσοι σπόροι περισσεύουν; Φτάνουν οι περισσευούμενοι για μία ακόμα γλάστρα;

\ex \textbf{Οι μπάλες.} Ο Άρης έχει 87 € και θέλει να αγοράσει μπάλες του μπάσκετ που κοστίζουν 13 € η καθεμία. Πόσες μπάλες αγοράζει και πόσα ευρώ του μένουν;

\ex \textbf{Η παράταξη.} Στην αυλή θα παραταχθούν 168 μαθητές της Α’ τάξης για μια φωτογραφία. Ο γυμναστής ρωτάει αν μπορούν να σταθούν σε ίσες ομάδες των 3, των 4, των 5, των 6 ή των 7, χωρίς να περισσεύει κανείς. Κάνε τις διαιρέσεις και απάντησε για κάθε περίπτωση.

\ex \textbf{Το ρολόι.} Η Μυρτώ βάζει το ξυπνητήρι στις 10:00 το πρωί. Θέλει να δει τι ώρα θα είναι μετά από 50 ώρες. Σκέψου ότι κάθε 24 ώρες το ρολόι γυρίζει στην ίδια ώρα.

\zoni{Το σκέφτομαι}{μια ερώτηση που πάει πιο μακριά}

\ex Βρες όλους τους φυσικούς μεταξύ του 20 και του 40 που, διαιρούμενοι με το 7, αφήνουν υπόλοιπο 3. Τι παρατηρείς για τη διαφορά ανάμεσα σε δύο διαδοχικούς από αυτούς;

\ex Ένας αριθμός $\nu$ αφήνει υπόλοιπο 4 όταν διαιρεθεί με το 6. Δοκίμασε μερικές τιμές του $\nu$ (10, 16, 22, \dots). Τι υπόλοιπο αφήνει ο $\nu$ όταν διαιρεθεί με το 3;

\ex Στη διαίρεση ενός αριθμού με έναν άλλον το πηλίκο είναι 12 και το υπόλοιπο 7. Ποιος είναι ο μικρότερος δυνατός διαιρέτης; Ποιος είναι τότε ο διαιρετέος;

\ex Η Ζωή ρωτάει: «Γιατί δεν διαιρούμε με το 0;» Ο Ηλίας απαντάει «γιατί δεν το επιτρέπει ο δάσκαλος». Εξήγησε με το $6:0=x$ και την αντίστροφη πράξη ότι δεν είναι θέμα άδειας. Τι γίνεται στο $0:0$;

\ex Σήμερα είναι Τρίτη. Τι μέρα θα είναι μετά από ένα εκατομμύριο ημέρες; (Μην αρχίσεις να μετράς. Βρες το υπόλοιπο της διαίρεσης του $1.000.000$ με το 7.)

\kefalaio{Χαρακτήρες διαιρετότητας: ξέρω χωρίς να διαιρέσω}{Κεφάλαιο Α.1.5, «Χαρακτήρες διαιρετότητας – ΜΚΔ – ΕΚΠ – Ανάλυση αριθμού σε γινόμενο πρώτων παραγόντων» (πρώτο μέρος)}

\begin{anagke}
Το σχολικό παζάρι πήγε καλύτερα απ' όσο περίμενε κανείς. Η Φανή, που κρατάει το ταμείο, μετράει τα χρήματα: 2.358 €. Θέλουν να τα μοιράσουν ίσα σε τρεις ομάδες που θα φτιάξουν από ένα σχολικό έργο. Μετά η Μυρτώ λέει: «Και αν οι ομάδες ήταν εννιά;» Ο Άρης: «Και αν ήταν πέντε;» Ο Στέλιος ξεκινάει να κάνει τις διαιρέσεις με τη σειρά, μία μία. Τρεις διαιρέσεις με τετραψήφιο αριθμό είναι πολλή δουλειά.

Η Λίνα κοιτάζει τον αριθμό 2.358 και λέει: «Για το πέντε δεν χρειάζεται καν να κάνεις διαίρεση. Τελειώνει σε οκτώ. Δεν γίνεται». «Πώς το ξέρεις;» ρωτάει ο Ηλίας. «Το έχω δει σε όλους τους αριθμούς που διαιρούνται με το πέντε: τελειώνουν σε 0 ή σε 5». Η Ζωή ρωτάει: «Και για το τρία και για το εννιά υπάρχει κάτι τέτοιο;»

Υπάρχει. Και είναι από τα πιο χρήσιμα κόλπα στο γυμνάσιο: να καταλαβαίνεις αν ένας αριθμός διαιρείται με έναν άλλον κοιτάζοντας απλώς τα ψηφία του.
\end{anagke}

\begin{lexi}{κριτήριο, χαρακτήρας διαιρετότητας}{κριτήριο, χαρακτήρας}
Το \lexiword{κριτήριο} έρχεται από το «κρίνω», δηλαδή αποφασίζω. Κριτήριο είναι ο κανόνας με τον οποίο κρίνουμε κάτι. Ένας κριτής σε αγώνα έχει κριτήρια για να βγάλει τον νικητή. Ένα \lexiword{κριτήριο διαιρετότητας} μας βοηθάει να κρίνουμε αν ένας αριθμός διαιρείται με κάποιον άλλον, χωρίς να κάνουμε τη διαίρεση.

Το σχολικό βιβλίο τα λέει και \lexiword{χαρακτήρες διαιρετότητας}. Ο χαρακτήρας είναι το «σημάδι» που χαρακτηρίζει κάτι, όπως ο χαρακτήρας ενός ανθρώπου. Εδώ το σημάδι είναι στα ψηφία του αριθμού.
\end{lexi}

\begin{metafrasi}
Ο αριθμός είναι άρτιος. & Λήγει σε $0,\,2,4,\,6$ ή $8$ (διαιρείται με το 2)\\
Ο αριθμός διαιρείται με το 5. & Λήγει σε $0$ ή $5$\\
Ο αριθμός διαιρείται με το 10. & Λήγει σε $0$\\
Ο αριθμός διαιρείται με το 3. & Το άθροισμα των ψηφίων του διαιρείται με το 3\\
Ο αριθμός διαιρείται με το 9. & Το άθροισμα των ψηφίων του διαιρείται με το 9\\
Ο αριθμός διαιρείται με το 4. & Ο αριθμός από τα δύο τελευταία ψηφία διαιρείται με το 4\\
Ο αριθμός διαιρείται με το 25. & Λήγει σε $00$, $25$, $50$ ή $75$\\
Ο αριθμός διαιρείται με το 6. & Διαιρείται και με το 2 και με το 3\\
\end{metafrasi}

\tmima{Η θεωρία}

\begin{kanonas}{κριτήρια διαιρετότητας}
\begin{itemize}[leftmargin=1.2em,itemsep=0.3ex,topsep=2pt]
\item \textbf{Με το 2}: το τελευταίο ψηφίο είναι 0, 2, 4, 6 ή 8.
\item \textbf{Με το 5}: το τελευταίο ψηφίο είναι 0 ή 5.
\item \textbf{Με το 10}: το τελευταίο ψηφίο είναι 0. (Με το 100: τα δύο τελευταία ψηφία είναι 00, με το 1.000: τα τρία τελευταία, κ.ο.κ.)
\item \textbf{Με το 4}: ο αριθμός που σχηματίζουν τα δύο τελευταία ψηφία διαιρείται με το 4 (ή είναι 00).
\item \textbf{Με το 25}: τα δύο τελευταία ψηφία είναι 00, 25, 50 ή 75.
\item \textbf{Με το 3}: το άθροισμα των ψηφίων διαιρείται με το 3.
\item \textbf{Με το 9}: το άθροισμα των ψηφίων διαιρείται με το 9.
\end{itemize}
\end{kanonas}

Από τα κριτήρια αυτά φτιάχνουμε και άλλα. Για να διαιρείται ένας αριθμός με το 6 πρέπει να διαιρείται και με το 2 και με το 3. Για το 15, και με το 3 και με το 5. Για το 12, και με το 3 και με το 4.

\tmimasmall{Γιατί δουλεύουν}

Κάθε αριθμός ξεκινάει με πολλαπλάσια του 10, του 100, του 1.000. Για παράδειγμα $3.725=37\cdot 100+25$. Ο 100 διαιρείται με το 4 και με το 25, άρα το $37\cdot 100$ διαιρείται και με τους δύο, και το μόνο που μένει να ελέγξουμε είναι τα δύο τελευταία ψηφία, το 25. Το ίδιο και για τα 2 και 5, με το 10.

Για το 3 και το 9 δουλεύει ένα κόλπο. Οι αριθμοί 9, 99, 999, \dots\ διαιρούνται με το 9 (και με το 3). Γράφω $2.457=2\cdot 1.000+4\cdot 100+5\cdot 10+7$ και αντικαθιστώ το $1.000=999+1$, $100=99+1$, $10=9+1$:
$$2.457=2\cdot 999+4\cdot 99+5\cdot 9+(2+4+5+7).$$
Τα τρία πρώτα κομμάτια διαιρούνται με το 9. Άρα όλος ο αριθμός διαιρείται με το 9 ακριβώς όταν διαιρείται το $2+4+5+7=18$. Και διαιρείται.

\begin{prosoxi}
\begin{itemize}[leftmargin=1.2em,itemsep=0.3ex,topsep=2pt]
\item Αν ένας αριθμός τελειώνει σε 4, δεν σημαίνει ότι διαιρείται με το 4. Το 14 δεν διαιρείται. Για το 4 κοιτάμε τα δύο τελευταία ψηφία μαζί.
\item Για το 3 και το 9 μετράνε \textbf{όλα} τα ψηφία, όχι μόνο το τελευταίο.
\item Για να διαιρείται ένας αριθμός με το 24 δεν φτάνει να διαιρείται με το 4 και με το 6. Το 12 διαιρείται και με τα δύο, αλλά όχι με το 24. Μπορούμε να «σπάσουμε» έναν αριθμό σε δύο μόνο όταν οι δύο αριθμοί δεν έχουν κοινό παράγοντα (όπως το 3 και το 4, ή το 2 και το 5).
\item Το 0 διαιρείται με κάθε αριθμό (εκτός από το 0 ως διαιρέτη), γιατί $0=\delta\cdot 0$.
\end{itemize}
\end{prosoxi}

\tmima{Λυμένα παραδείγματα}

\paradeigma{Ο πίνακας των κριτηρίων}
\tixero Οι αριθμοί 12.510, 772, 225 και 13.600.\par
\tithelo Με ποιους από τους 2, 3, 4, 5, 9, 10, 25 και 100 διαιρείται ο καθένας.\par
\possindeo Εφαρμόζω κάθε κριτήριο. Στα 3 και 9 υπολογίζω το άθροισμα ψηφίων. Στα 4 και 25 κοιτάζω τα δύο τελευταία ψηφία.

\begin{center}
\small\renewcommand{\arraystretch}{1.3}\setlength{\tabcolsep}{4pt}
\begin{tabular}{@{}lcccccccc@{}}
\toprule
 & 2 & 3 & 4 & 5 & 9 & 10 & 25 & 100\\
\midrule
12.510 \ ($1+2+5+1+0=9$) & ναι & ναι & όχι & ναι & ναι & ναι & όχι & όχι\\
772 \ ($7+7+2=16$, και $72:4=18$) & ναι & όχι & ναι & όχι & όχι & όχι & όχι & όχι\\
225 \ ($2+2+5=9$) & όχι & ναι & όχι & ναι & ναι & όχι & ναι & όχι\\
13.600 \ ($1+3+6=10$) & ναι & όχι & ναι & ναι & όχι & ναι & ναι & ναι\\
\bottomrule
\end{tabular}
\end{center}

\apantisi Δες τον πίνακα.\par
\elegxos Το 13.600 διαιρείται με το 100, άρα και με το 10, το 4, το 25, το 5 και το 2, όπως πρέπει: κάθε διαιρέτης του 100 μπαίνει αυτόματα. Δοκιμή με το 4: $13.600:4=3.400$. Σωστά.

\paradeigma{Το ψηφίο που λείπει}
\tixero Ο αριθμός $4.7\square 2$, όπου το $\square$ είναι ένα ψηφίο.\par
\tithelo Ποιο ψηφίο πρέπει να μπει ώστε ο αριθμός να διαιρείται με το 9.\par
\possindeo Το άθροισμα των γνωστών ψηφίων είναι $4+7+2=13$. Πρέπει $13+\square$ να διαιρείται με το 9. Τα πολλαπλάσια του 9 κοντά στο 13 είναι το 18, άρα $\square=5$. (Το επόμενο πολλαπλάσιο, 27, θα απαιτούσε ψηφίο 14, που δεν υπάρχει.)\par
\apantisi $4.752$.\par
\elegxos $4.752:9=528$, γιατί $9\cdot 528=4.752$. Σωστά.

\paradeigma{Η μοιρασιά του παζαριού}
\tixero Η Φανή έχει 2.358 €.\par
\tithelo Σε ποιες από τις 3, 9, 5, 4 και 6 ισόποσες ομάδες μοιράζονται χωρίς να περισσεύει τίποτα.\par
\possindeo Άθροισμα ψηφίων: $2+3+5+8=18$. Διαιρείται με το 3 και με το 9. Τελευταίο ψηφίο 8, άρα όχι με το 5. Τελευταία δύο ψηφία 58: το $58:4$ δεν είναι ακέραιο, άρα όχι με το 4. Είναι άρτιος και διαιρείται με το 3, άρα διαιρείται με το 6.\par
\apantisi Σε 3, 6 και 9 ομάδες.\par
\elegxos $2.358:3=786$, $2.358:6=393$, $2.358:9=262$. Και οι τρεις είναι τέλειες διαιρέσεις.

\tmima{Ασκήσεις}

\zoni{Το κατάλαβα}{εφαρμογή του κανόνα}

\ex Ποιοι από τους αριθμούς 3.140, 2.025, 1.234, 5.050 και 707 διαιρούνται με το 2; Ποιοι με το 5; Ποιοι με το 10;

\ex Βρες ποιοι από τους 123, 4.521, 1.001, 99.999 και 2.468 διαιρούνται με το 3. Έλεγξε μία από τις περιπτώσεις με διαίρεση.

\ex Βρες ποιοι από τους 1.089, 4.527, 81.081, 7.777 και 90.009 διαιρούνται με το 9.

\ex Βρες ποιοι από τους 316, 1.002, 7.500, 2.024 και 98 διαιρούνται με το 4.

\ex Βρες ποιοι από τους 175, 2.350, 4.600, 1.065 και 3.025 διαιρούνται με το 25.

\ex Συμπλήρωσε έναν πίνακα με τους αριθμούς 840, 1.575, 6.048 και 20.250 και τους διαιρέτες 2, 3, 4, 5, 9, 10, 25. Για κάθε αριθμό γράψε ναι ή όχι.

\ex Ποιοι από τους 54, 128, 4.230, 1.005 και 7.776 διαιρούνται με το 6;

\ex Βρες το ψηφίο (ή τα ψηφία) που πρέπει να μπουν στη θέση του $\square$: α) $4.7\square 2$ να διαιρείται με το 9 \quad β) $3\square 5$ να διαιρείται με το 3 \quad γ) $6\square 4$ να διαιρείται με το 4 \quad δ) $17\square$ να διαιρείται και με το 2 και με το 5.

\ex Σωστό ή λάθος; α) Κάθε αριθμός που τελειώνει σε 4 διαιρείται με το 4. \quad β) Αν ένας αριθμός διαιρείται με το 9, διαιρείται και με το 3. \quad γ) Το 0 διαιρείται με το 7. \quad δ) Αν ένας αριθμός διαιρείται με το 4 και με το 6, διαιρείται και με το 24. \quad ε) Αν διαιρείται με το 2 και με το 5, διαιρείται με το 10. \quad στ) Αν διαιρείται με το 3 και με το 4, διαιρείται με το 12.

\zoni{Το χρησιμοποιώ}{προβλήματα από τη ζωή}

\ex \textbf{Τα χρήματα του παζαριού.} Δες το παράδειγμα της Φανής, αλλά τώρα τα χρήματα είναι 3.465 €. Σε ποιες από τις 3, 5, 9, 4 και 10 ομάδες μοιράζονται ακριβώς; Βρες πόσα παίρνει η κάθε ομάδα, όπου γίνεται.

\ex \textbf{Τα μπλουζάκια.} Ο προπονητής του Άρη έχει 140 μπλουζάκια και θέλει να τα βάλει σε ίσες στοίβες. Με ποιους από τους αριθμούς 2, 3, 4, 5, 6 και 10 γίνεται, χωρίς να περισσεύει κανένα; Χρησιμοποίησε κριτήρια.

\ex \textbf{Ο κωδικός της Λίνας.} Ο νέος κωδικός της Λίνας είναι τετραψήφιος. Ξεκινάει με 7, το δεύτερο ψηφίο είναι το 2, τελειώνει σε 0, διαιρείται με το 9 και το τρίτο ψηφίο είναι μεγαλύτερο από το δεύτερο. Ποιος είναι ο κωδικός;

\ex \textbf{Οι θέσεις του λεωφορείου.} Οι θέσεις ενός λεωφορείου είναι αριθμημένες από το 1 ως το 60. Οι θέσεις που διαιρούνται με το 5 είναι δίπλα στο παράθυρο. Πόσες είναι; Πόσες διαιρούνται με το 10; Πόσες διαιρούνται και με το 3 και με το 5;

\ex \textbf{Τα δίσεκτα έτη.} Ένα έτος είναι δίσεκτο αν διαιρείται με το 4. Εξαίρεση οι αιώνες (1.900, 2.000, \dots), που είναι δίσεκτοι μόνο όταν διαιρούνται με το 400. Ποια από τα έτη 1.900, 2.000, 2.026, 2.028 και 2.100 είναι δίσεκτα;

\zoni{Το σκέφτομαι}{μια ερώτηση που πάει πιο μακριά}

\ex Γιατί ένας αριθμός που τελειώνει σε 0 διαιρείται και με το 2 και με το 5; (Βοήθεια: $10=2\cdot 5$.)

\ex Γράψε το 1.452 ως $1\cdot 999+4\cdot 99+5\cdot 9+(1+4+5+2)$ και εξήγησε με τη βοήθεια αυτής της γραφής αν διαιρείται με το 3 και αν διαιρείται με το 9.

\ex Ποιος είναι ο μικρότερος τετραψήφιος αριθμός που διαιρείται ταυτόχρονα με τα 2, 3, 4, 5 και 9;

\ex Οι αριθμοί 121, 132, 143 και 154 διαιρούνται όλοι με το 11. Τι παρατηρείς στα ψηφία τους; Δοκίμασε το ίδιο μοτίβο στο 165 και στο 175 και έλεγξε με διαίρεση.

\ex Ο Στέλιος ισχυρίζεται ότι ο 2.026 διαιρείται με το 3, γιατί τελειώνει σε 6, και το 6 διαιρείται με το 3. Ποιο είναι το λάθος του; Τι θα έπρεπε να είχε κάνει;

\kefalaio{Πρώτοι αριθμοί και ανάλυση: τα δομικά υλικά}{Κεφάλαιο Α.1.5, «Χαρακτήρες διαιρετότητας – ΜΚΔ – ΕΚΠ – Ανάλυση αριθμού σε γινόμενο πρώτων παραγόντων» (δεύτερο μέρος)}

\begin{anagke}
Ο Θοδωρής, ο μεγάλος αδελφός του Άρη, έχει περάσει να τους δει μετά το σχολείο. Στο λύκειο έχει αρχίσει να παίζει με κρυπτογραφία, το πώς φυλάσσονται οι κωδικοί στο διαδίκτυο, και δεν χάνει ευκαιρία να το πει. «Σκέφτομαι δύο πρώτους αριθμούς», λέει στη Λίνα, που τον κοιτάει με ενδιαφέρον. «Το γινόμενό τους είναι 143. Ποιοι είναι;»

Η Λίνα δοκιμάζει. Δεν διαιρείται με το 2, γιατί τελειώνει σε 3. Το άθροισμα ψηφίων $1+4+3=8$, άρα ούτε με το 3. Ούτε με το 5. Με το 7, $7\cdot 20=140$ και περισσεύουν 3, όχι. Με το 11: $11\cdot 13=143$. «Έντεκα και δεκατρία!» φωνάζει. «Μπράβο», λέει ο Θοδωρής. «Τώρα σκέψου ότι ο αριθμός έχει 600 ψηφία. Ένας υπολογιστής πολλαπλασιάζει τους δύο πρώτους σε ένα κλάσμα του δευτερολέπτου. Για να τους βρει από το γινόμενο, θέλει απίστευτα πολύ χρόνο. Πάνω σε αυτό στηρίζονται οι κωδικοί των τραπεζών».

Η Ζωή ρωτάει κάτι πιο απλό. «Δηλαδή κάθε αριθμός είναι φτιαγμένος από πρώτους, όπως ο τοίχος από τούβλα;» Ναι. Και το να βρίσκεις αυτά τα τούβλα είναι δουλειά με μέθοδο.
\end{anagke}

\begin{lexi}{πρώτος, σύνθετος, παράγοντας}{πρώτος, σύνθετος}
Ο \lexiword{πρώτος} αριθμός δεν λέγεται έτσι επειδή έρχεται πρώτος στη σειρά, αλλά επειδή είναι «πρωταρχικός», πρωτογενής: δεν φτιάχνεται από μικρότερους αριθμούς με πολλαπλασιασμό. Είναι το δομικό υλικό.

Ο \lexiword{σύνθετος} αριθμός είναι αυτός που συντίθεται, δηλαδή φτιάχνεται, από άλλους: το 12 είναι $2\cdot 2\cdot 3$. Ο \lexiword{παράγοντας} είναι ο αριθμός που παίρνει μέρος σε έναν πολλαπλασιασμό και «παράγει» το γινόμενο.
\end{lexi}

\begin{metafrasi}
Το 7 έχει διαιρέτες μόνο το 1 και το 7. & Το 7 είναι πρώτος\\
Το 12 έχει διαιρέτες τους 1, 2, 3, 4, 6, 12. & Το 12 είναι σύνθετος\\
Το 60 σαν γινόμενο πρώτων. & $60=2\cdot 2\cdot 3\cdot 5=2^2\cdot 3\cdot 5$\\
Ο $\alpha$ διαιρεί τον $\beta$. & $\beta=\alpha\cdot\kappa$, με $\kappa$ φυσικό\\
Πολλαπλάσια του 4. & $0,\,4,\,8,\,12,\,16,\,20,\,\dots$\\
\end{metafrasi}

\tmima{Η θεωρία}

\tmimasmall{Πολλαπλάσια και διαιρέτες}

\textbf{Πολλαπλάσια} ενός φυσικού $\alpha$ είναι οι αριθμοί που βγαίνουν αν τον πολλαπλασιάσουμε με φυσικούς: $0,\ \alpha,\ 2\alpha,\ 3\alpha,\ 4\alpha,\ \dots$ \textbf{Διαιρέτες} του $\alpha$ είναι όλοι οι αριθμοί που τον διαιρούν ακριβώς. Κάθε αριθμός $\alpha$ έχει διαιρέτες τουλάχιστον τον 1 και τον εαυτό του. Αν ο $\delta$ διαιρεί έναν αριθμό, τότε ο αριθμός είναι πολλαπλάσιο του $\delta$. Είναι το ίδιο πράγμα, από δύο πλευρές.

\begin{orismos}{πρώτος και σύνθετος}
\textbf{Πρώτος} λέγεται ένας φυσικός αριθμός μεγαλύτερος από το 1 που έχει διαιρέτες μόνο τον 1 και τον εαυτό του. Ένας φυσικός μεγαλύτερος από το 1 που έχει και άλλους διαιρέτες λέγεται \textbf{σύνθετος}. Το 1 δεν είναι ούτε πρώτος ούτε σύνθετος.
\end{orismos}

Οι πρώτοι μέχρι το 50 είναι οι: \blue{2, 3, 5, 7, 11, 13, 17, 19, 23, 29, 31, 37, 41, 43, 47}. Ο 2 είναι ο μοναδικός άρτιος πρώτος, γιατί κάθε άλλος άρτιος διαιρείται και με το 2. Οι πρώτοι είναι άπειροι: το απέδειξε ο Ευκλείδης πριν από 2.300 χρόνια.

\tmimasmall{Το κόσκινο του Ερατοσθένη}

Ο Ερατοσθένης, που έζησε στην Αλεξάνδρεια τον 3ο αιώνα π.Χ., βρήκε μια μέθοδο για να «κοσκινίζουμε» τους πρώτους. Γράφεις τους αριθμούς, σβήνεις το 1, μετά κρατάς το 2 και σβήνεις όλα τα πολλαπλάσιά του, μετά κρατάς το επόμενο που δεν έσβησες, το 3, και σβήνεις τα πολλαπλάσιά του, και ούτω καθεξής. Αυτά που μένουν είναι οι πρώτοι.

\begin{center}
\renewcommand{\arraystretch}{1.25}
\setlength{\tabcolsep}{5pt}
\newcommand{\pr}[1]{\textcolor{gblue}{\bfseries #1}}
\begin{tabular}{cccccccccc}
1 & \pr{2} & \pr{3} & 4 & \pr{5} & 6 & \pr{7} & 8 & 9 & 10\\
\pr{11} & 12 & \pr{13} & 14 & 15 & 16 & \pr{17} & 18 & \pr{19} & 20\\
21 & 22 & \pr{23} & 24 & 25 & 26 & 27 & 28 & \pr{29} & 30\\
\pr{31} & 32 & 33 & 34 & 35 & 36 & \pr{37} & 38 & 39 & 40\\
\pr{41} & 42 & \pr{43} & 44 & 45 & 46 & \pr{47} & 48 & 49 & 50\\
\end{tabular}
\end{center}

\tmimasmall{Ανάλυση σε γινόμενο πρώτων παραγόντων}

Κάθε σύνθετος αριθμός γράφεται ως γινόμενο πρώτων, και \textbf{με έναν μόνο τρόπο} (αν αγνοήσουμε τη σειρά των παραγόντων). Γι' αυτό οι πρώτοι είναι τα δομικά υλικά: κάθε αριθμός έχει το δικό του «σύνολο υλικών».

\begin{kanonas}{ανάλυση σε πρώτους παράγοντες}
\begin{enumerate}[leftmargin=1.6em,itemsep=0.2ex,topsep=2pt]
\item Βρίσκω τον μικρότερο πρώτο που διαιρεί τον αριθμό (δοκιμάζω με τη σειρά 2, 3, 5, 7, 11, \dots με τα κριτήρια διαιρετότητας).
\item Διαιρώ και γράφω τον πρώτο στα δεξιά και το πηλίκο από κάτω.
\item Επαναλαμβάνω με το πηλίκο, ώσπου να φτάσω στο 1.
\item Γράφω το γινόμενο των πρώτων και τους ίδιους παράγοντες τους ενώνω σε δυνάμεις.
\end{enumerate}
\end{kanonas}

Αν ο αριθμός είναι πρώτος, δεν χρειάζεται να δοκιμάσεις όλους τους πρώτους. Δοκιμάζεις με τη σειρά και σταματάς όταν το \textbf{πηλίκο γίνει μικρότερο από τον διαιρέτη}. Αν μέχρι τότε δεν βρήκες διαιρέτη, ο αριθμός είναι πρώτος.

\begin{prosoxi}
\begin{itemize}[leftmargin=1.2em,itemsep=0.3ex,topsep=2pt]
\item Το 1 \textbf{δεν είναι πρώτος}. Αν ήταν, το 6 θα γραφόταν και $2\cdot 3$ και $1\cdot 2\cdot 3$ και $1\cdot 1\cdot 2\cdot 3$, και η ανάλυση δεν θα ήταν μοναδική.
\item Ο 2 είναι πρώτος, και μάλιστα ο μόνος άρτιος. Όχι «όλοι οι πρώτοι είναι περιττοί».
\item Ο αριθμός που μοιάζει πρώτος δεν είναι πάντα. Το 91 φαίνεται πρώτος, αλλά $91=7\cdot 13$.
\end{itemize}
\end{prosoxi}

\tmima{Λυμένα παραδείγματα}

\paradeigma{Πρώτος ή σύνθετος;}
\tixero Οι αριθμοί 91 και 97.\par
\tithelo Να αποφασίσω αν είναι πρώτοι.\par
\possindeo Για το 91: δεν διαιρείται με το 2 (τελειώνει σε 1), ούτε με το 3 ($9+1=10$), ούτε με το 5. Με το 7: $7\cdot 13=91$. Βρήκα διαιρέτη, άρα είναι σύνθετος.\par
Για το 97: ούτε με 2, ούτε με 3 ($9+7=16$), ούτε με 5. Με το 7: $97=7\cdot 13+6$, δεν διαιρείται, και το πηλίκο 13 είναι ακόμα μεγαλύτερο από το 7. Με το 11: $97=11\cdot 8+9$, και το πηλίκο 8 είναι μικρότερο από το 11. Σταματάω.\par
\apantisi Το 91 είναι σύνθετος ($7\cdot 13$), το 97 είναι πρώτος.\par
\elegxos Οι πρώτοι κάτω από το 100 είναι 25 και ο τελευταίος από αυτούς είναι το 97.

\paradeigma{Ανάλυση του 2.520}
\tixero Ο αριθμός 2.520.\par
\tithelo Να τον αναλύσω σε γινόμενο πρώτων παραγόντων.\par
\possindeo Διαιρώ διαδοχικά: $2.520:2=1.260$, $1.260:2=630$, $630:2=315$. Το 315 είναι περιττός, άρα δοκιμάζω με το 3: $315:3=105$, $105:3=35$. Μετά με το 5: $35:5=7$. Τέλος $7:7=1$.\par
Οι παράγοντες είναι τρία δυάρια, δύο τριάρια, ένα πεντάρι και ένα εφτάρι.\par
\apantisi $2.520=2^3\cdot 3^2\cdot 5\cdot 7$.\par
\elegxos $8\cdot 9=72$, $72\cdot 5=360$, $360\cdot 7=2.520$. Σωστά.

\paradeigma{Το μυστικό του Θοδωρή}
\tixero Το 143 είναι το γινόμενο δύο πρώτων.\par
\tithelo Ποιοι είναι αυτοί οι πρώτοι.\par
\possindeo Δοκιμάζω τους πρώτους με τη σειρά. Το 2, το 3, το 5 και το 7 δεν διαιρούν το 143. Με το 11: $143:11=13$, και το 13 είναι πρώτος.\par
\apantisi $143=11\cdot 13$.\par
\elegxos $11\cdot 13=143$. Και οι δύο παράγοντες είναι πρώτοι, άρα αυτή είναι η ανάλυση.

\tmima{Ασκήσεις}

\zoni{Το κατάλαβα}{εφαρμογή του κανόνα}

\ex Γράψε όλους τους διαιρέτες των αριθμών 12, 18, 30, 36 και 45.

\ex Γράψε τα έξι πρώτα πολλαπλάσια (μαζί με το 0) των αριθμών 4, 7 και 12.

\ex Ποιοι από τους 1, 2, 9, 13, 21, 29, 33, 41, 51, 57, 61, 87 και 89 είναι πρώτοι; Για τους σύνθετους, γράψε έναν διαιρέτη τους.

\ex Βρες όλους τους πρώτους αριθμούς ανάμεσα στο 50 και στο 80.

\ex Ανάλυσε σε γινόμενο πρώτων παραγόντων (με δυνάμεις): α) 24 \quad β) 60 \quad γ) 72 \quad δ) 100 \quad ε) 150 \quad στ) 504.

\ex Ποιοι αριθμοί έχουν ανάλυση: α) $2^3\cdot 5$ \quad β) $2\cdot 3^2\cdot 7$ \quad γ) $2^2\cdot 3\cdot 5^2$ \quad δ) $3^3\cdot 11$;

\ex Σωστό ή λάθος; α) Κάθε πρώτος αριθμός είναι περιττός. \quad β) Το 1 είναι πρώτος. \quad γ) Το 91 είναι πρώτος. \quad δ) Υπάρχουν άπειροι πρώτοι αριθμοί. \quad ε) Το 0 είναι πολλαπλάσιο κάθε αριθμού. \quad στ) Το 51 είναι πρώτος.

\zoni{Το χρησιμοποιώ}{προβλήματα από τη ζωή}

\ex \textbf{Το μυστικό του Θοδωρή.} Ο Θοδωρής σκέφτεται δύο πρώτους με γινόμενο 221. Ποιοι είναι; Δοκίμασε με τη σειρά τους πρώτους μέχρι να βρεις διαιρέτη.

\ex \textbf{Οι καραμέλες της Μυρτώς.} Η Μυρτώ έχει 48 καραμέλες και θέλει να τις βάλει σε ίσα σακουλάκια, χωρίς να περισσέψει καμία. Πόσες καραμέλες μπορεί να βάλει σε κάθε σακουλάκι; Η Μυρτώ έχει και ένα άλλο βάζο με 47 καραμέλες. Με πόσους τρόπους μπορεί να τις μοιράσει εκεί; Γιατί τόσο λίγους;

\ex \textbf{Οι ομάδες.} Ο γυμναστής θέλει να χωρίσει 60 παιδιά σε ίσες ομάδες. Θέλει τουλάχιστον 2 ομάδες και τουλάχιστον 2 παιδιά σε κάθε ομάδα. Με πόσους διαφορετικούς τρόπους γίνεται; Αν τα παιδιά ήταν 61, πόσοι τρόποι θα υπήρχαν;

\ex \textbf{Η φανέλα του Άρη.} Ο αριθμός στη φανέλα του Άρη είναι πρώτος αριθμός ανάμεσα στο 20 και στο 30, και το άθροισμα των ψηφίων του είναι 11. Ποιος είναι ο αριθμός;

\ex \textbf{Το κόσκινο.} Φτιάξε το κόσκινο του Ερατοσθένη για τους αριθμούς από το 1 ως το 60. Πόσοι πρώτοι υπάρχουν; Πόσους αριθμούς σβήνει το 2, και πόσους ακόμα το 3;

\zoni{Το σκέφτομαι}{μια ερώτηση που πάει πιο μακριά}

\ex Βρες τους διαιρέτες του 6 και του 28, εκτός από τον ίδιο τον αριθμό, και πρόσθεσέ τους. Τι παρατηρείς; Οι μαθηματικοί λένε αυτούς τους αριθμούς «τέλειους».

\ex Δύο πρώτοι που διαφέρουν κατά 2 λέγονται δίδυμοι, όπως οι 3 και 5. Βρες όλα τα ζευγάρια διδύμων πρώτων μέχρι το 50.

\ex Ο 72 είναι $2^3\cdot 3^2$. Γράψε όλους τους αριθμούς της μορφής $2^{\alpha}\cdot 3^{\beta}$, όπου $\alpha=0,\,1,2,\,3$ και $\beta=0,\,1,2$. Τι είναι αυτοί οι αριθμοί για τον 72; Πόσοι είναι;

\ex Η Ζωή ρωτάει «γιατί το 1 δεν είναι πρώτος, αφού διαιρείται μόνο με τον εαυτό του και με το 1;» Εξήγησε με την ανάλυση του 6 γιατί, αν το 1 ήταν πρώτος, δεν θα είχε κάθε αριθμός μία μοναδική ανάλυση.

\ex Ο Θοδωρής ισχυρίζεται ότι $10.403=101\cdot 103$. Έλεγξέ το με πολλαπλασιασμό. Πόσο δύσκολο θα ήταν να βρεις τους δύο παράγοντες χωρίς να τους ξέρεις, δοκιμάζοντας τους πρώτους με τη σειρά;

\kefalaio{ΜΚΔ και ΕΚΠ: μοιράζω και συναντιέμαι}{Κεφάλαιο Α.1.5, «Χαρακτήρες διαιρετότητας – ΜΚΔ – ΕΚΠ – Ανάλυση αριθμού σε γινόμενο πρώτων παραγόντων» (τρίτο μέρος)}

\begin{anagke}
Αυτή την εβδομάδα η κυρία Αλεξάνδρα βάζει δύο προβλήματα που μοιάζουν να έχουν αντίθετη ιδέα.

Πρώτο: το σχολείο φτιάχνει κήπο. Η Α’ τάξη έχει 150 φυτά, 90 γλάστρες και 60 σακούλες χώμα. Θέλουν να φτιάξουν ίδια κιτ για τα τμήματα, ώστε κάθε κιτ να έχει και από τα τρία είδη και να μη μείνει τίποτα. «Πόσα το πολύ ίδια κιτ γίνονται;» Ο Στέλιος λέει αμέσως «δέκα». Η Φανή ελέγχει: δέκα κιτ θέλουν 15 φυτά, 9 γλάστρες και 6 σακούλες το καθένα. Βγαίνει. «Μα δεν είπαμε πόσα το πολύ;» λέει η Ζωή.

Δεύτερο: ο θείος Λευτέρης και ο Ηλίας κάθονται το βράδυ στην ακτή και βλέπουν δύο φάρους. Ο ένας αναβοσβήνει κάθε 12 δευτερόλεπτα, ο άλλος κάθε 18. Μόλις άναψαν και οι δύο μαζί. «Σε πόσο ξανά θα ανάψουν μαζί;» ρωτάει ο θείος.

Στο πρώτο θέλουμε το μεγαλύτερο κοινό κομμάτι: μοιράζουμε. Στο δεύτερο θέλουμε την πρώτη φορά που συμπίπτουν: συναντιόμαστε. Και τα δύο λύνονται με τους ίδιους αριθμούς, διαιρέτες και πολλαπλάσια.
\end{anagke}

\begin{lexi}{ΜΚΔ και ΕΚΠ}{μέγιστος, ελάχιστος, κοινός}
Οι δύο συντομογραφίες κρύβουν τρεις λέξεις η καθεμία. \lexiword{Κοινός} είναι αυτός που ανήκει και στους δύο. \lexiword{Μέγιστος} είναι ο μεγαλύτερος που υπάρχει (το «μέγας» στον υπερθετικό), και \lexiword{ελάχιστος} ο μικρότερος που υπάρχει (το «ελαχύς»).

Το \lexiword{ΜΚΔ} είναι ο Μέγιστος Κοινός Διαιρέτης: ο μεγαλύτερος αριθμός που διαιρεί και τους δύο. Το \lexiword{ΕΚΠ} είναι το Ελάχιστο Κοινό Πολλαπλάσιο: το μικρότερο (μη μηδενικό) πολλαπλάσιο που είναι και των δύο.
\end{lexi}

\begin{metafrasi}
Ο μεγαλύτερος αριθμός που διαιρεί και το 12 και το 18. & $\MKD(12,\,18)=6$\\
Ο μικρότερος αριθμός (όχι 0) που είναι πολλαπλάσιο και του 12 και του 18. & $\EKP(12,\,18)=36$\\
Οι $\alpha$ και $\beta$ δεν έχουν κοινό διαιρέτη εκτός από το 1. & $\MKD(\alpha,\beta)=1$ (πρώτοι μεταξύ τους)\\
Ο $\alpha$ διαιρεί τον $\beta$. & $\MKD(\alpha,\beta)=\alpha$, $\EKP(\alpha,\beta)=\beta$\\
Το γινόμενο των δύο αριθμών. & $\MKD\cdot\EKP=\alpha\cdot\beta$\\
\end{metafrasi}

\tmima{Η θεωρία}

\begin{orismos}{ΜΚΔ}
Ο \textbf{Μέγιστος Κοινός Διαιρέτης} δύο ή περισσότερων φυσικών είναι ο μεγαλύτερος από τους κοινούς τους διαιρέτες. Γράφουμε $\MKD(\alpha,\beta)$.
\end{orismos}

\begin{orismos}{ΕΚΠ}
Το \textbf{Ελάχιστο Κοινό Πολλαπλάσιο} δύο ή περισσότερων φυσικών είναι το μικρότερο από τα κοινά τους πολλαπλάσια, εκτός από το 0 (που είναι πολλαπλάσιο όλων). Γράφουμε $\EKP(\alpha,\beta)$.
\end{orismos}

\tmimasmall{Δύο μέθοδοι}

\textbf{Με λίστες.} Γράφω τους διαιρέτες (για ΜΚΔ) ή τα πρώτα πολλαπλάσια (για ΕΚΠ) κάθε αριθμού και διαλέγω τον κοινό που ψάχνω. Είναι εύκολο για μικρούς αριθμούς.

\begin{kanonas}{ΜΚΔ και ΕΚΠ με ανάλυση}
\begin{enumerate}[leftmargin=1.6em,itemsep=0.2ex,topsep=2pt]
\item Αναλύω κάθε αριθμό σε γινόμενο πρώτων παραγόντων (με δυνάμεις).
\item \textbf{ΜΚΔ}: παίρνω μόνο τους \emph{κοινούς} πρώτους παράγοντες, καθένα με τον \emph{μικρότερο} εκθέτη που έχει.
\item \textbf{ΕΚΠ}: παίρνω \emph{όλους} τους πρώτους παράγοντες (κοινούς και μη), καθένα με τον \emph{μεγαλύτερο} εκθέτη που έχει.
\item Πολλαπλασιάζω τους παράγοντες που πήρα.
\end{enumerate}
\end{kanonas}

Η ιδέα: το ΜΚΔ πρέπει να χωράει και στους δύο αριθμούς, γι' αυτό παίρνει όσα έχουν και οι δύο, το λιγότερο. Το ΕΚΠ πρέπει να χωράει και τους δύο, γι' αυτό παίρνει όλα όσα χρειάζεται ο καθένας, το περισσότερο.

\tmimasmall{Πρώτοι μεταξύ τους και μια σχέση}

Δύο αριθμοί λέγονται \textbf{πρώτοι μεταξύ τους} αν ο μοναδικός κοινός διαιρέτης τους είναι το 1, δηλαδή αν $\MKD(\alpha,\beta)=1$. Δεν χρειάζεται να είναι οι ίδιοι πρώτοι: οι 8 και 15 είναι σύνθετοι αλλά πρώτοι μεταξύ τους. Τότε το ΕΚΠ τους είναι το γινόμενό τους: $\EKP(8,\,15)=120$.

Για δύο φυσικούς ισχύει πάντα
$$\MKD(\alpha,\beta)\cdot\EKP(\alpha,\beta)=\alpha\cdot\beta.$$
Για τα 12 και 18: $6\cdot 36=216=12\cdot 18$.

\begin{prosoxi}
\begin{itemize}[leftmargin=1.2em,itemsep=0.3ex,topsep=2pt]
\item Το ΜΚΔ είναι πάντα μικρότερο ή ίσο από τον μικρότερο αριθμό. Το ΕΚΠ είναι πάντα μεγαλύτερο ή ίσο από τον μεγαλύτερο. Αν βγάλεις ΜΚΔ μεγαλύτερο από έναν από τους αριθμούς, κάτι πήγε στραβά.
\item Στο ΜΚΔ προσοχή: παίρνω μόνο τους \emph{κοινούς} παράγοντες. Στο ΕΚΠ παίρνω \emph{όλους}.
\item Όταν ένας αριθμός διαιρεί τον άλλον, τα πράγματα είναι εύκολα: $\MKD(7,42)=7$ και $\EKP(7,42)=42$.
\end{itemize}
\end{prosoxi}

\tmima{Λυμένα παραδείγματα}

\paradeigma{Οι φάροι}
\tixero Ο ένας φάρος αναβοσβήνει κάθε 12 δευτερόλεπτα, ο άλλος κάθε 18. Άναψαν μαζί τώρα.\par
\tithelo Σε πόσα δευτερόλεπτα θα ανάψουν ξανά μαζί.\par
\possindeo Ο πρώτος ανάβει στα δευτερόλεπτα $12, 24, 36, 48, \dots$ και ο δεύτερος στα $18, 36, 54, \dots$ Η πρώτη φορά που συμπίπτουν είναι το 36. Δηλαδή ψάχνω το $\EKP(12,18)$.\par
Με ανάλυση: $12=2^2\cdot 3$ και $18=2\cdot 3^2$. Παίρνω όλους τους παράγοντες με τον μεγαλύτερο εκθέτη: $2^2\cdot 3^2=36$.\par
\apantisi Σε 36 δευτερόλεπτα.\par
\elegxos Το 36 είναι πολλαπλάσιο και του 12 ($12\cdot 3$) και του 18 ($18\cdot 2$), και κανένας μικρότερος αριθμός δεν είναι (δες τις λίστες).

\paradeigma{Τα κιτ του κήπου}
\tixero 150 φυτά, 90 γλάστρες και 60 σακούλες χώμα. Τα κιτ πρέπει να είναι ίδια και να μη μείνει τίποτα.\par
\tithelo Πόσα το πολύ κιτ γίνονται και τι έχει το καθένα.\par
\possindeo Ο αριθμός των κιτ πρέπει να διαιρεί και το 150 και το 90 και το 60. Ψάχνω το μεγαλύτερο κοινό διαιρέτη. $150=2\cdot 3\cdot 5^2$, $90=2\cdot 3^2\cdot 5$, $60=2^2\cdot 3\cdot 5$. Κοινοί παράγοντες: 2, 3 και 5, με τον μικρότερο εκθέτη τον 1 σε κάθε περίπτωση. Άρα $\MKD=2\cdot 3\cdot 5=30$.\par
Κάθε κιτ έχει: $150:30=5$ φυτά, $90:30=3$ γλάστρες, $60:30=2$ σακούλες.\par
\apantisi 30 κιτ, με 5 φυτά, 3 γλάστρες και 2 σακούλες το καθένα.\par
\elegxos Τα 10 κιτ του Στέλιου βγαίνουν, αλλά το 30 είναι κοινός διαιρέτης και μεγαλύτερος από το 10. Αν δοκιμάσουμε μεγαλύτερο αριθμό, π.χ. 60 κιτ, το 90 δεν διαιρείται με το 60 ($90:60$ δεν είναι ακέραιο).

\paradeigma{ΜΚΔ και ΕΚΠ τριών αριθμών}
\tixero Οι αριθμοί 2.520, 2.940 και 3.780.\par
\tithelo Το ΜΚΔ και το ΕΚΠ τους.\par
\possindeo Από το προηγούμενο κεφάλαιο: $2.520=2^3\cdot 3^2\cdot 5\cdot 7$. Με τον ίδιο τρόπο $2.940=2^2\cdot 3\cdot 5\cdot 7^2$ και $3.780=2^2\cdot 3^3\cdot 5\cdot 7$.\par
ΜΚΔ: οι παράγοντες 2, 3, 5, 7 υπάρχουν και στους τρεις. Παίρνω τον μικρότερο εκθέτη του καθενός: $2^2\cdot 3\cdot 5\cdot 7=420$.\par
ΕΚΠ: παίρνω τον μεγαλύτερο εκθέτη του καθενός: $2^3\cdot 3^3\cdot 5\cdot 7^2=8\cdot 27\cdot 5\cdot 49=52.920$.\par
\apantisi $\MKD=420$ και $\EKP=52.920$.\par
\elegxos Το 420 διαιρεί τον καθένα: $2.520:420=6$, $2.940:420=7$, $3.780:420=9$. Και το 52.920 είναι πολλαπλάσιο του καθενός: $52.920:2.520=21$.

\tmima{Ασκήσεις}

\zoni{Το κατάλαβα}{εφαρμογή του κανόνα}

\ex Βρες το ΜΚΔ γράφοντας τους διαιρέτες: α) 12 και 30 \quad β) 18 και 24 \quad γ) 20 και 45 \quad δ) 16 και 27.

\ex Βρες το ΕΚΠ γράφοντας τα πρώτα πολλαπλάσια: α) 4 και 6 \quad β) 6 και 15 \quad γ) 8 και 12 \quad δ) 5 και 7.

\ex Βρες το ΜΚΔ και το ΕΚΠ με ανάλυση σε πρώτους παράγοντες: α) 24 και 36 \quad β) 60 και 84 \quad γ) 90 και 150 \quad δ) 48 και 180.

\ex Βρες το ΜΚΔ και το ΕΚΠ των τριών αριθμών 12, 18 και 30.

\ex Ποια από τα ζευγάρια είναι πρώτοι μεταξύ τους; α) 8 και 15 \quad β) 12 και 35 \quad γ) 21 και 35 \quad δ) 9 και 28 \quad ε) 14 και 49 \quad στ) 1 και 99.

\ex Βρες το ΜΚΔ και το ΕΚΠ των: α) 7 και 42 \quad β) 15 και 60 \quad γ) 25 και 100. Τι παρατηρείς όταν ο ένας αριθμός διαιρεί τον άλλον;

\ex Σωστό ή λάθος; α) $\MKD(12,18)=36$. \quad β) $\EKP(4,6)=24$. \quad γ) Το ΜΚΔ δύο διαφορετικών πρώτων είναι 1. \quad δ) Το ΕΚΠ δύο αριθμών δεν είναι ποτέ μικρότερο από τον μεγαλύτερο. \quad ε) Το ΜΚΔ δεν είναι ποτέ μεγαλύτερο από τον μικρότερο. \quad στ) Αν δύο αριθμοί είναι πρώτοι μεταξύ τους, το ΕΚΠ τους είναι το γινόμενό τους.

\zoni{Το χρησιμοποιώ}{προβλήματα από τη ζωή}

\ex \textbf{Οι φάροι.} Δύο άλλοι φάροι αναβοσβήνουν ο ένας κάθε 15 δευτερόλεπτα και ο άλλος κάθε 20. Άναψαν μαζί τώρα. Πότε θα ξανανάψουν μαζί; Πόσες φορές θα συμπέσουν στα επόμενα 3 λεπτά, μετρώντας και το τέλος του τρίτου λεπτού;

\ex \textbf{Τα καλάθια της γιαγιάς.} Η γιαγιά Ρήνη έχει 84 μήλα και 126 πορτοκάλια και θέλει να φτιάξει ίδια καλάθια, με και από τα δύο φρούτα και χωρίς να περισσέψει τίποτα. Πόσα το πολύ καλάθια μπορεί να φτιάξει; Τι έχει το καθένα;

\ex \textbf{Τα πλακάκια του θείου.} Η βεράντα του θείου Λευτέρη είναι ορθογώνια, 180 εκ. επί 150 εκ. Θέλει να την καλύψει με τετράγωνα πλακάκια ίδιου μεγέθους, χωρίς να κόψει κανένα. Ποια είναι η μεγαλύτερη δυνατή πλευρά του πλακιδίου; Πόσα πλακάκια θα χρειαστεί;

\ex \textbf{Τα λεωφορεία.} Δύο γραμμές λεωφορείων ξεκινούν μαζί από τη στάση στις 7:00. Η μία περνάει κάθε 15 λεπτά και η άλλη κάθε 20. Πότε θα ξανακαταλήξουν μαζί; Πόσες φορές θα βρεθούν μαζί από τις 7:00 ως τις 10:00 (μετρώντας και τις δύο ώρες);

\ex \textbf{Το πότισμα.} Η Μυρτώ ποτίζει τις γλάστρες κάθε 4 μέρες, ο Ηλίας κάθε 6 και η Ζωή κάθε 9. Σήμερα τις πότισαν και οι τρεις. Σε πόσες μέρες θα τις ποτίσουν ξανά και οι τρεις την ίδια μέρα;

\zoni{Το σκέφτομαι}{μια ερώτηση που πάει πιο μακριά}

\ex Για τα ζευγάρια 12 και 18, 8 και 20, 9 και 15, 14 και 21 βρες το ΜΚΔ, το ΕΚΠ και το γινόμενο των δύο αριθμών. Σύγκρινε το γινόμενο ΜΚΔ·ΕΚΠ με το γινόμενο των αριθμών. Τι ανακαλύπτεις;

\ex Η Ζωή παρατηρεί ότι $48=18\cdot 2+12$, και ότι $\MKD(48,18)=6=\MKD(18,12)$. Δοκίμασε το ίδιο με τα 84 και 30: γράψε την Ευκλείδεια διαίρεση, και σύγκρινε το $\MKD(84,30)$ με το $\MKD(30,\upsilon)$. Αν συνεχίσεις μέχρι να βρεις υπόλοιπο 0, τι δείχνει ο τελευταίος διαιρέτης; Δοκίμασε με τα 252 και 105. (Αυτή είναι η μέθοδος του Ευκλείδη.)

\ex Γιατί, όταν δύο αριθμοί είναι πρώτοι μεταξύ τους, το ΕΚΠ είναι το γινόμενό τους; Εξήγησε με τα 8 και 15, χρησιμοποιώντας τις αναλύσεις τους.

\ex Βρες δύο αριθμούς που να έχουν ΜΚΔ 6 και ΕΚΠ 36. Υπάρχει και δεύτερο ζευγάρι; (Βοήθεια: το γινόμενο των δύο πρέπει να είναι $6\cdot 36$.)

\ex Μπορεί το ΜΚΔ δύο \emph{διαφορετικών} αριθμών να είναι ίσο με το ΕΚΠ τους; Εξήγησε.

\epanalipsi{1}

\tmima{Λεξικό της ενότητας}

\renewcommand{\arraystretch}{1.5}
\begin{tabular*}{\linewidth}{@{\extracolsep{\fill}} >{\raggedright\arraybackslash}p{0.26\linewidth} >{\raggedright\arraybackslash}p{0.64\linewidth}}
\toprule
\textbf{Λέξη ή τύπος} & \textbf{Με μία πρόταση}\\
\midrule
φυσικός & Ένας από τους $0,\,1,2,\,3,\dots$: μετράει πράγματα που δεν σπάνε.\\
στρογγυλοποίηση & Αλλάζω έναν αριθμό με κοντινό του, με συμφωνημένα βήματα, μία φορά.\\
$\alpha^{\nu}$ & Ο $\alpha$ γραμμένος $\nu$ φορές σε γινόμενο.\\
προτεραιότητα & Παρενθέσεις, δυνάμεις, $\cdot$ και $:$, και τέλος $+$ και $-$.\\
$\Delta=\delta\cdot\pi+\upsilon$ & Ευκλείδεια διαίρεση, με $\upsilon<\delta$.\\
κριτήρια & Κανόνες στα ψηφία που δείχνουν αν ένας αριθμός διαιρείται με το 2, 3, 4, 5, 9, 10, 25.\\
πρώτος & Ο αριθμός μεγαλύτερος από το 1 με διαιρέτες μόνο τον 1 και τον εαυτό του.\\
ΜΚΔ & Ο μεγαλύτερος κοινός διαιρέτης: μοιράζω.\\
ΕΚΠ & Το μικρότερο (μη μηδενικό) κοινό πολλαπλάσιο: συναντιέμαι.\\
\bottomrule
\end{tabular*}

\tmima{Ασκήσεις επανάληψης}

\ex Γράψε με ψηφία τον αριθμό «τριάντα επτά εκατομμύρια, δύο χιλιάδες, πέντε». Μετά στρογγυλοποίησέ τον στα εκατομμύρια.

\ex Βάλε σε αύξουσα σειρά τους 5.012, 5.102, 5.021, 4.999 και 10.000.

\ex Στρογγυλοποίησε το 86.549 στις δεκάδες, στις εκατοντάδες και στις χιλιάδες.

\ex Υπολόγισε με τις ιδιότητες των πράξεων: α) $25\cdot 17\cdot 4$ \quad β) $99\cdot 8$ \quad γ) $36\cdot 15+36\cdot 85$ \quad δ) $47+68+53+32$.

\ex Βρες την τιμή: α) $3\cdot(8-2)^2-10$ \quad β) $2^5-5^2$ \quad γ) $(2+3)\cdot 4-2^3$ \quad δ) $100-4\cdot 5^2$.

\ex Κάνε την Ευκλείδεια διαίρεση $1.000:48$ και γράψε την ισότητα $\Delta=\delta\cdot\pi+\upsilon$. Ποιος αριθμός δίνει πηλίκο 7 και υπόλοιπο 6 όταν διαιρεθεί με το 15;

\ex Σήμερα είναι Πέμπτη. Τι μέρα θα είναι μετά από 100 ημέρες;

\ex Με ποιους από τους 2, 3, 4, 5, 9, 10 και 25 διαιρείται ο αριθμός 3.240;

\ex Ποιο ψηφίο πρέπει να μπει στη θέση του $\square$ ώστε ο $5\square 7$ να διαιρείται με το 9;

\ex Ποιοι από τους 57, 59, 63 και 67 είναι πρώτοι; Για τους σύνθετους, βρες την ανάλυσή τους.

\ex Ανάλυσε τον 360 σε γινόμενο πρώτων παραγόντων. Πόσους πρώτους διαφορετικούς παράγοντες έχει;

\ex Βρες το ΜΚΔ και το ΕΚΠ των αριθμών 36 και 60.

\ex \textbf{Οι κάρτες του Στέλιου.} Ο Στέλιος έχει 36 κάρτες με ζώα και 60 κάρτες με αυτοκίνητα. Θέλει να φτιάξει όσο το δυνατόν περισσότερες ίδιες στοίβες, με και από τα δύο είδη, χωρίς να περισσέψει καμία. Πόσες στοίβες θα γίνουν και τι έχει η κάθε μία;

\ex \textbf{Η πισίνα.} Ο Ηλίας πάει στην πισίνα κάθε 8 μέρες και η Ζωή κάθε 12. Σήμερα πήγαν μαζί. Πότε θα ξαναπάνε μαζί;

\ex Σωστό ή λάθος; α) Κάθε άρτιος αριθμός διαιρείται με το 4. \quad β) $\MKD(12,\,18)=6$. \quad γ) Το 91 είναι πρώτος. \quad δ) $2^4=8$. \quad ε) $17-(9-3)=17-9-3$. \quad στ) $0:5=0$.

\ex Ο Άρης λέει ότι $6+4\cdot 5=50$ και η Ζωή ότι $6+4\cdot 5=26$. Ποιος έχει δίκιο και γιατί; Πού πρέπει να βάλεις παρένθεση ώστε να βγει 50;

\ex Πόσους διαιρέτες έχει το 100; Πόσοι από αυτούς είναι πρώτοι;

\tmima{Ο χάρτης της επανάληψης}

Μετά τις λύσεις, δες σε ποιες ασκήσεις δυσκολεύτηκες. Ο πίνακας σου λέει ποιο κεφάλαιο να ξαναδιαβάσεις.

\renewcommand{\arraystretch}{1.5}
\begin{tabular*}{\linewidth}{@{\extracolsep{\fill}} >{\raggedright\arraybackslash}p{0.38\linewidth} >{\raggedright\arraybackslash}p{0.52\linewidth}}
\toprule
\textbf{Αν δυσκολεύτηκες στις} & \textbf{ξαναδιάβασε τα κεφάλαια}\\
\midrule
Ε1.1, Ε1.2, Ε1.3 & 1 (φυσικοί, διάταξη, στρογγυλοποίηση)\\
Ε1.4 & 2 (ιδιότητες των πράξεων)\\
Ε1.5, Ε1.16 & 3 (δυνάμεις, προτεραιότητα)\\
Ε1.6, Ε1.7 & 4 (Ευκλείδεια διαίρεση)\\
Ε1.8, Ε1.9 & 5 (κριτήρια διαιρετότητας)\\
Ε1.10, Ε1.11, Ε1.17 & 6 (πρώτοι, ανάλυση)\\
Ε1.12, Ε1.13, Ε1.14 & 7 (ΜΚΔ, ΕΚΠ)\\
Ε1.15 & 1 έως 7\\
\bottomrule
\end{tabular*}



\backmatter
\cleardoublepage
\addcontentsline{toc}{part}{Λύσεις}
\markboth{Λύσεις}{Λύσεις}
{\color{gblue}\bfseries\Large Λύσεις}\par\vspace{2mm}
{\color{ggray}\small Κοίτα τες αφού προσπαθήσεις, όχι πριν.\par}
\lyseis{1}{Φυσικοί αριθμοί: η σειρά που δεν τελειώνει}
\sol{1} α) 205 \quad β) 732 \quad γ) 20.813 \quad δ) 4.500.015.
\sol{2} α) τριάντα οκτώ χιλιάδες εννιακόσια πενήντα ένα \quad β) πέντε εκατομμύρια οκτακόσια δώδεκα \quad γ) εκατόν είκοσι χιλιάδες τρία.
\sol{3} Οι τρεις προηγούμενοι του 289 είναι οι 288, 287, 286 και οι δύο επόμενοι οι 290, 291. Ο προηγούμενος του 1.000.000 είναι ο 999.999. Το 0 δεν έχει προηγούμενο στους φυσικούς.
\sol{4} Αύξουσα: 3.508, 3.515, 3.620, 4.800, 4.801. Φθίνουσα: 4.801, 4.800, 3.620, 3.515, 3.508.
\sol{5} α) $=$ \quad β) $>$ \quad γ) $<$ \quad δ) $<$ \quad ε) $>$ (9 > 8 στο πρώτο ψηφίο) \quad στ) $<$ (λιγότερα ψηφία).
\sol{6} Κάθε 2 εκ. είναι μία μονάδα. Β: $6:2=3$, Γ: $10:2=5$, Δ: $12:2=6$, Ε: $14:2=7$.
\sol{7} $345\approx 300$ (δεκάδες 4), $761\approx 800$ (6), $659\approx 700$ (5), $2.567\approx 2.600$ (6), $9.532\approx 9.500$ (3), $31.549\approx 31.500$ (4), $8.950\approx 9.000$ (5, και το 9 των εκατοντάδων «κουβαλάει»).
\sol{8} α) $7.568.350$ (μονάδες 9) \quad β) $7.568.300$ (δεκάδες 4) \quad γ) $7.568.000$ (εκατοντάδες 3) \quad δ) $7.570.000$ (χιλιάδες 8, άρα το 6 γίνεται 7) \quad ε) $8.000.000$ (εκατοντάδες χιλιάδων 5, άρα το 7 γίνεται 8).
\sol{9} α) Σωστό. β) Σωστό. γ) Σωστό: $122-32+1=91$. δ) Σωστό: $35-12+1=24$. ε) Σωστό: οκτώ ημέρες είναι μία εβδομάδα και μία μέρα. στ) Σωστό. ζ) Λάθος: το ψηφίο των δεκάδων είναι 4, άρα $1.449\approx 1.400$.
\sol{10} 306, 360, 603, 630. Ο 036 θα ήταν ο αριθμός 36, με δύο ψηφία, γιατί το 0 μπροστά δεν μετράει.
\sol{11} Η Λίνα στρογγυλοποίησε στις χιλιάδες ($2.917\approx 3.000$) και ο Ηλίας στις εκατοντάδες ($2.917\approx 2.900$). Το 2.900 απέχει 17 από το 2.917, το 3.000 απέχει 83. Πιο κοντά είναι ο Ηλίας.
\sol{12} Όχι: $9.987<10.000$. Του λείπουν $10.000-9.987=13$ βήματα. Ο στόχος λέει «τουλάχιστον», δηλαδή χρειάζεται ακρίβεια. Η στρογγυλοποίηση έκρυψε τα 13 βήματα που έλειπαν.
\sol{13} Στις εκατοντάδες $1.234\approx 1.200$, αλλά 12 κιβώτια έχουν 1.200 πλακάκια και λείπουν 34. Χρειάζονται 13 κιβώτια, δηλαδή 1.300 πλακάκια, και περισσεύουν $1.300-1.234=66$. Εδώ στρογγυλοποιούμε προς τα πάνω, ό,τι κι αν λέει ο κανόνας.
\sol{14} $112-47+1=66$ σελίδες. Σε 10 ημέρες διαβάζει $10\times 6=60$ σελίδες, άρα δεν φτάνουν: μένουν 6. Χρειάζονται 11 ημέρες, γιατί $66:6=11$.
\sol{15} Τα 20 μ. χωρίζονται σε $20:2=10$ διαστήματα. Οι πάσσαλοι είναι ένας παραπάνω από τα διαστήματα: 11 πάσσαλοι.
\sol{16} $399-120+1=280$ εισιτήρια. Ξεκινούν από άρτιο (120) και τελειώνουν σε περιττό (399), άρα εναλλάσσονται άρτιος, περιττός και τα μισά είναι περιττά: 140.
\sol{17} Τα 100 αντιστοιχούν σε 2 εκ., άρα κάθε εκατοστό δείχνει 50. Α: $3\times 50=150$, Β: $5\times 50=250$, Γ: $7{,}5\times 50=375$, Δ: $10\times 50=500$.
\sol{18} Στις μονάδες το 9 βρίσκεται στους 9, 19, 29, \dots, 99: 10 φορές. Στις δεκάδες βρίσκεται στους 90 ως 99: 10 φορές. Σύνολο 20 φορές (το 99 έχει δύο εννιάρια).
\sol{19} Ο μικρότερος είναι ο 2.500 (εκατοντάδες 5, άρα προς τα πάνω) και ο μεγαλύτερος ο 3.499. Ο 2.499 δίνει 2.000 και ο 3.500 δίνει 4.000.
\sol{20} Ο μεγαλύτερος είναι ο 9.876 και ο μικρότερος ο 1.023. Δεν είναι ο 0.123, γιατί δεν είναι τετραψήφιος.
\sol{21} Αν υπήρχε μεγαλύτερος φυσικός, έστω Μ, τότε ο επόμενός του $Μ+1$ θα ήταν φυσικός και μεγαλύτερος από το Μ. Αυτό είναι αδύνατο, άρα δεν υπάρχει.
\sol{22} Όχι: ο $3-5$ δεν είναι φυσικός, γιατί δεν μπορείς να αφαιρέσεις 5 από 3. Θα χρειαστούμε αριθμούς «κάτω από το 0», τους αρνητικούς, που θα γνωρίσεις στο Τεύχος Β.

\lyseis{2}{Πρόσθεση, αφαίρεση, πολλαπλασιασμός: οι ιδιότητες που κρύβονται}
\sol{1} Στο $13+5=18$: προσθετέοι και άθροισμα. Στο $20-8=12$: μειωτέος, αφαιρετέος και διαφορά. Στο $7\cdot 6=42$: παράγοντες και γινόμενο.
\sol{2} α) αντιμεταθετική της πρόσθεσης \quad β) προσεταιριστική της πρόσθεσης \quad γ) επιμεριστική του πολλαπλασιασμού ως προς την πρόσθεση \quad δ) το 1 δεν αλλάζει το γινόμενο \quad ε) ο πολλαπλασιασμός με το 0 δίνει 0 \quad στ) το 0 δεν αλλάζει το άθροισμα.
\sol{3} α) 5.200 \quad β) 370 \quad γ) 4.900.000 \quad δ) 8.000.
\sol{4} α) $40-17=23$ \quad β) $90-26=64$ \quad γ) $32+15=47$ \quad δ) $80-35=45$. Έλεγχος, π.χ. στο δ: $80-45=35$.
\sol{5} α) $4\cdot 25\cdot 17=100\cdot 17=1.700$ \quad β) $8\cdot 125\cdot 39=1.000\cdot 39=39.000$ \quad γ) $(12+88)+(47+53)=100+100=200$ \quad δ) $5\cdot 2\cdot 73=10\cdot 73=730$.
\sol{6} α) $89\cdot(7+3)=89\cdot 10=890$ \quad β) $(23+77)\cdot 49=100\cdot 49=4.900$ \quad γ) $76\cdot(13-3)=76\cdot 10=760$ \quad δ) $284\cdot(100-1)=28.400-284=28.116$ \quad ε) $35\cdot(100+2)=3.500+70=3.570$.
\sol{7} α) Λάθος: το $4-9$ δεν υπάρχει στους φυσικούς. β) Λάθος: $(20-8)-2=10$, ενώ $20-(8-2)=14$. γ) Λάθος: $5\cdot(2+3)=25$, ενώ $5\cdot 2+3=13$. δ) Λάθος: $7\cdot 0=0$. ε) Σωστό. στ) Σωστό.
\sol{8} α) 190 \quad β) $(122+78)+25=225$ \quad γ) 462 \quad δ) 2.726 (έλεγχος: $4.595+2.726=7.321$) \quad ε) $60-16=44$.
\sol{9} $1+2+3+4=10$, \ $1+2+34=37$, \ $12+34=46$, \ $1\cdot 2\cdot 3\cdot 4=24$.
\sol{10} $98\cdot 5=(100-2)\cdot 5=500-10=490$ τοστ.
\sol{11} Αφετηρία ως γέφυρα: $47$ χλμ. Γέφυρα ως βενζινάδικο: $112-47=65$ χλμ. Βενζινάδικο ως χωριό: $175-112=63$ χλμ. Χωριό ως θάλασσα: $230-175=55$ χλμ. Βενζινάδικο ως θάλασσα: $230-112=118$ χλμ. Έλεγχος: $63+55=118$ και $47+65+63+55=230$.
\sol{12} Μειωτέος είναι τα αυτοκόλλητα που είχε: $Μ=17+38=55$. Έλεγχος: $55-17=38$.
\sol{13} Στην αφαίρεση $82-47$ το 2 στις μονάδες είναι μικρότερο από το 7. Δεν επιτρέπεται να γυρίσουμε τους αριθμούς ανάποδα: δανειζόμαστε μία δεκάδα και κάνουμε $12-7=5$, και μετά $7-4=3$. Σωστό: 35. Έλεγχος με $Μ=Α+\Delta$: $47+35=82$. Με το 45 του Άρη βγαίνει $47+45=92$, που δεν είναι 82.
\sol{14} $48\cdot 25=12\cdot 4\cdot 25=12\cdot 100=1.200$ πλακάκια.
\sol{15} $13\cdot 8=(10+3)\cdot 8=10\cdot 8+3\cdot 8=80+24=104$ τ.μ. Ένα ορθογώνιο $10\times 8$ και ένα $3\times 8$, κολλημένα.
\sol{16} $3\cdot 99=3\cdot(100-1)=300-3=297$ €.
\sol{17} Τα ζευγάρια είναι 50 και το καθένα κάνει 101, άρα $50\cdot 101=5.050$. Για το $1+\dots+10$: 5 ζευγάρια $(1+10),(2+9),\dots$ και καθένα κάνει 11, άρα $5\cdot 11=55$. Έλεγχος με πρόσθεση: $1+2+3+4+5+6+7+8+9+10=55$.
\sol{18} Οι αριθμοί είναι 7 και 5: $7+5=12$ και $7-5=2$.
\sol{19} Τέσσερις σακούλες με 0 καραμέλες έχουν συνολικά $0+0+0+0=0$ καραμέλες. Γενικά $4\cdot 0=0+0+0+0=0$.
\sol{20} Το ορθογώνιο με 3 σειρές των 5 έχει 15 τετράγωνα. Αν το γυρίσεις στο πλάι, βλέπεις 5 σειρές των 3, και τα τετράγωνα είναι πάλι 15. Η εικόνα είναι διαφορετική (3 ομάδες των 5 ή 5 ομάδες των 3), το αποτέλεσμα ίδιο.
\sol{21} $(10-5)-2=5-2=3$ και $10-(5-2)=10-3=7$. Οι τιμές διαφέρουν, άρα η αφαίρεση δεν είναι προσεταιριστική: έχει σημασία ποιες δύο αφαιρέσεις γίνονται πρώτα.
\sol{22} $3\cdot 3=9$, $9\cdot 3=27$, $27\cdot 3=81$, $81\cdot 3=243$, $243\cdot 3=729$. Έχει 3 ψηφία. Πιο σύντομα: $3^6$, «το τρία στην έκτη».

\lyseis{3}{Δυνάμεις και προτεραιότητα: ποιος πάει πρώτος}
\sol{1} α) $5^3$ \quad β) $2^5$ \quad γ) $10^4$ \quad δ) $7^2$ \quad ε) $9=9^1$ (ή $3^2$).
\sol{2} α) $2\cdot 2\cdot 2\cdot 2=16$ \quad β) $3\cdot 3\cdot 3=27$ \quad γ) $5\cdot 5=25$ \quad δ) $100.000$ \quad ε) $1$ \quad στ) $6$.
\sol{3} Στο $4^5$: βάση 4, εκθέτης 5. Στο $12^2$: βάση 12, εκθέτης 2. Στο $1^9$: βάση 1, εκθέτης 9. Στο $10^8$: βάση 10, εκθέτης 8.
\sol{4} Για τους 2, 3, 4, 5, 10: τετράγωνα 4, 9, 16, 25, 100 και κύβοι 8, 27, 64, 125, 1.000.
\sol{5} α) $12+12=24$ \quad β) $15\cdot 4=60$ \quad γ) $24:4=6$, μετά $6:2=3$ \quad δ) $100-20=80$ \quad ε) $2+9=11$ \quad στ) $5^2=25$.
\sol{6} α) $10.100$ (δες το παράδειγμα). β) $5^3-8\cdot 9=125-72=53$. γ) $3\cdot 6^2-10\cdot 3=3\cdot 36-30=108-30=78$.
\sol{7} α) $7\cdot 10^3+6\cdot 10^2+0\cdot 10+4$ \quad β) $5\cdot 10^4+3\cdot 10^3+2\cdot 10^2+8\cdot 10$ \quad γ) $4\cdot 10^5+6\cdot 10+5$.
\sol{8} α) $3.205$ \quad β) $900.040$ \quad γ) $1.000.100$.
\sol{9} α) Λάθος: $2^3=8$. β) Σωστό. γ) Λάθος: $5^2=25$, ενώ $2\cdot 5=10$. δ) Σωστό. ε) Λάθος: πρώτα ο πολλαπλασιασμός, $3+10=13$. στ) Λάθος: $(1+1)^2=4$, ενώ $1^2+1^2=2$.
\sol{10} Με 3 τροχούς: $10^3=1.000$ κωδικοί. Με 4: $10^4=10.000$. Με 5: $10^5=100.000$. Κάθε τροχός που προστίθεται πολλαπλασιάζει τους κωδικούς επί 10.
\sol{11} Με 3 διακόπτες: $2^3=8$ συνδυασμοί. Με 5: $2^5=32$. Με 10: $2^{10}=1.024$. Κάθε διακόπτης που προστίθεται διπλασιάζει τους συνδυασμούς.
\sol{12} $4^3=4\cdot 4\cdot 4=64$ κυβικά δεκατόμετρα, δηλαδή 64 λίτρα.
\sol{13} Εμβαδό $12^2=144$ τ.μ. Με πλευρά 24 μ. το εμβαδό είναι $24^2=576$ τ.μ., που είναι $4\cdot 144$. Όχι διπλάσιο, τετραπλάσιο.
\sol{14} $8^2=64$ τετράγωνα.
\sol{15} Ταιριάζει η $3\cdot 8+12$: τρεις μπάλες των 8 € και ένα φανελάκι των 12 €. Πληρώνει $24+12=36$ €. Η $3\cdot(8+12)=60$ € θα ταίριαζε αν αγόραζε τρία ζευγάρια «μπάλα και φανελάκι».
\sol{16} Σωστή είναι η Ζωή (300). Στην παράσταση $4\cdot(7+7\cdot 9)+20$ πρώτα γίνεται ο πολλαπλασιασμός μέσα στην παρένθεση: $7\cdot 9=63$, μετά $7+63=70$, μετά $4\cdot 70=280$ και $280+20=300$. Ο Ηλίας (524) έκανε πρώτα το $7+7=14$ μέσα στην παρένθεση: $4\cdot 14\cdot 9+20=524$. Ο Άρης (335) έκανε πρώτα το $4\cdot 7+7=35$ και μετά επί 9: $35\cdot 9+20=335$. Και οι δύο παραβίασαν την προτεραιότητα.
\sol{17} $2^3\cdot 2^2=8\cdot 4=32$. Ως παράγοντες: $(2\cdot 2\cdot 2)\cdot(2\cdot 2)$, δηλαδή πέντε δυάρια: $2^5=32$. Ο κανόνας: προσθέτω τους εκθέτες, $2^3\cdot 2^2=2^{3+2}$.
\sol{18} $2^5=32>5^2=25$ και $2^4=16=4^2$. Ο $\alpha=0$ κάνει το $\alpha^2=0$ ίσο με το $2\alpha=0$. (Μόνο το 0 και το 2 ικανοποιούν την ισότητα.)
\sol{19} Χωρίς παρενθέσεις: $2+12+5=19$. Με παρενθέσεις: $(2+3)\cdot 4+5=25$, $2+3\cdot(4+5)=2+27=29$, $(2+3)\cdot(4+5)=45$. Οι διαφορετικές τιμές είναι 19, 25, 29 και 45.
\sol{20} Π.χ. $1+2+3+4=10$ και $1\cdot 2\cdot 3+4=10$.
\sol{21} Το μοτίβο: $10^3=1.000$, $10^2=100$, $10^1=10$. Κάθε φορά που κατεβαίνει ο εκθέτης κατά 1, ο αριθμός διαιρείται με το 10. Το επόμενο είναι $10:10=1$, άρα $10^0=1$. Δίκιο έχει η Ζωή.

\lyseis{4}{Ευκλείδεια διαίρεση: το υπόλοιπο που μένει}
\sol{1} α) $4.002:69=58$ (δοκιμή: $69\cdot 58=4.002$) \quad β) $1.445:17=85$ \quad γ) $925:37=25$ \quad δ) $3.621:213=17$ \quad ε) $35.280:2.940=12$ \quad στ) $5.082:77=66$. Όλες είναι τέλειες διαιρέσεις, με υπόλοιπο 0.
\sol{2} α) $43=7\cdot 6+1$ \quad β) $100=9\cdot 11+1$ \quad γ) $245=12\cdot 20+5$ \quad δ) $1.000=33\cdot 30+10$ \quad ε) $59=60\cdot 0+59$ (το 60 δεν χωράει στο 59, πηλίκο 0) \quad στ) $0=5\cdot 0+0$.
\sol{3} α) Ναι, με διαιρέτη το 35 ($\upsilon=20<35$, πηλίκο 3). Όχι με διαιρέτη το 3, γιατί $20>3$. β) Όχι: το υπόλοιπο 40 είναι μεγαλύτερο και από το 38 και από το 19. γ) Ναι, με διαιρέτη το 42 (πηλίκο 35) και με διαιρέτη το 35 (πηλίκο 42), γιατί $30<35$. δ) Ναι, με διαιρέτη το 18 (πηλίκο 16) και με διαιρέτη το 16 (πηλίκο 18), γιατί $12<16$.
\sol{4} Τα υπόλοιπα είναι μικρότερα από το 8: 0, 1, 2, 3, 4, 5, 6 ή 7.
\sol{5} $\Delta=9\cdot 73+4=657+4=661$. Δοκιμή: $661:9$ δίνει πηλίκο 73 και υπόλοιπο 4.
\sol{6} α) $\Delta=12\cdot 7+5=89$ \quad β) $100=8\cdot 12+4$, άρα $\pi=12$, $\upsilon=4$ \quad γ) $59=\delta\cdot 4+3$, άρα $4\delta=56$, $\delta=14$. Έλεγχος: $3<14$.
\sol{7} α) Λάθος: το υπόλοιπο είναι μικρότερο από τον διαιρέτη, άρα το πολύ 5. β) Σωστό. γ) Λάθος: δεν διαιρούμε με το 0. δ) Σωστό. ε) Λάθος: $17:1=17$. στ) Σωστό.
\sol{8} Η ισότητα ισχύει: $12\cdot 6+13=72+13=85$. Όμως το υπόλοιπο 13 δεν είναι μικρότερο από τον διαιρέτη 12: ο 12 χωράει άλλη μία φορά στο 13. Σωστό: $85=12\cdot 7+1$.
\sol{9} $387=49\cdot 7+44$: χρειάζονται 8 λεωφορεία. Στο όγδοο υπάρχουν 44 παιδιά και μένουν $49-44=5$ κενές θέσεις. Με 7 λεωφορεία μένουν στο σχολείο 44 παιδιά.
\sol{10} $150=7\cdot 21+3$: Τρίτη και 3 μέρες, δηλαδή Παρασκευή. Πριν από 8 ημέρες: $8=7+1$, άρα μία μέρα πριν από την Τρίτη, δηλαδή Δευτέρα.
\sol{11} $1.000=33\cdot 30+10$: γεμίζει 30 γλάστρες και περισσεύουν 10 σπόροι. Όχι, δεν φτάνουν για άλλη γλάστρα, γιατί $10<33$.
\sol{12} $87=13\cdot 6+9$: αγοράζει 6 μπάλες (78 €) και του μένουν 9 €.
\sol{13} $168:3=56$, $168:4=42$, $168:6=28$, $168:7=24$: τέλειες διαιρέσεις, μπορούν να σταθούν σε ίσες ομάδες. Το $168:5$ δίνει πηλίκο 33 και υπόλοιπο 3 ($168=5\cdot 33+3$): περισσεύουν 3 παιδιά.
\sol{14} $50=24\cdot 2+2$: μετά από 2 ολόκληρα εικοσιτετράωρα το ρολόι δείχνει πάλι 10:00 και μένουν 2 ώρες. Θα είναι 12:00, δύο μέρες αργότερα.
\sol{15} Οι αριθμοί είναι 24, 31, 38 ($24=7\cdot 3+3$, $31=7\cdot 4+3$, $38=7\cdot 5+3$). Η διαφορά ανάμεσα σε δύο διαδοχικούς είναι πάντα 7, δηλαδή ο διαιρέτης.
\sol{16} Για $\nu=10$: $10:3$ έχει υπόλοιπο 1. Για 16: $16=3\cdot 5+1$, υπόλοιπο 1. Για 22: $22=3\cdot 7+1$, υπόλοιπο 1. Πάντα 1. Ο λόγος: $\nu=6\cdot\kappa+4$, το $6\kappa$ διαιρείται με το 3, και το 4 αφήνει υπόλοιπο 1.
\sol{17} Ο διαιρέτης πρέπει να είναι μεγαλύτερος από το υπόλοιπο 7, άρα ο μικρότερος είναι ο 8. Τότε $\Delta=8\cdot 12+7=96+7=103$.
\sol{18} Αν $6:0=x$, τότε $6=0\cdot x$, που είναι αδύνατο, γιατί το $0\cdot x$ είναι πάντα 0. Δεν υπάρχει $x$. Στο $0:0=x$ θα έπρεπε $0=0\cdot x$, που ισχύει για κάθε $x$: δεν υπάρχει ένας μοναδικός αριθμός. Και στις δύο περιπτώσεις η πράξη δεν ορίζεται.
\sol{19} $1.000.000=7\cdot 142.857+1$. Το υπόλοιπο είναι 1, άρα μία μέρα μετά την Τρίτη: Τετάρτη.

\lyseis{5}{Χαρακτήρες διαιρετότητας: ξέρω χωρίς να διαιρέσω}
\sol{1} Με το 2: 3.140, 1.234, 5.050. Με το 5: 3.140, 2.025, 5.050. Με το 10: 3.140, 5.050.
\sol{2} Άθροισμα ψηφίων: 123 $\to$ 6 (ναι), 4.521 $\to$ 12 (ναι), 1.001 $\to$ 2 (όχι), 99.999 $\to$ 45 (ναι), 2.468 $\to$ 20 (όχι). Έλεγχος, π.χ. $4.521:3=1.507$.
\sol{3} 1.089 $\to$ 18 (ναι), 4.527 $\to$ 18 (ναι), 81.081 $\to$ 18 (ναι), 7.777 $\to$ 28 (όχι), 90.009 $\to$ 18 (ναι).
\sol{4} 316 (16, ναι), 1.002 (02, όχι), 7.500 (00, ναι), 2.024 (24, ναι), 98 (98 δεν διαιρείται με το 4, όχι).
\sol{5} 175 (75, ναι), 2.350 (50, ναι), 4.600 (00, ναι), 1.065 (65, όχι), 3.025 (25, ναι).
\sol{6} \begin{tabular}[t]{@{}lccccccc@{}}\toprule & 2 & 3 & 4 & 5 & 9 & 10 & 25\\\midrule 840 & ναι & ναι & ναι & ναι & όχι & ναι & όχι\\ 1.575 & όχι & ναι & όχι & ναι & ναι & όχι & ναι\\ 6.048 & ναι & ναι & ναι & όχι & ναι & όχι & όχι\\ 20.250 & ναι & ναι & όχι & ναι & ναι & ναι & ναι\\\bottomrule\end{tabular}
\sol{7} Το 54 (άρτιος και $5+4=9$), το 4.230 (άρτιος και $4+2+3+0=9$) και το 7.776 (άρτιος και $7+7+7+6=27$). Το 128 δεν διαιρείται με το 3, και το 1.005 είναι περιττός.
\sol{8} α) $4+7+2=13$, χρειάζεται το 18, άρα $\square=5$: 4.752. \quad β) $3+5=8$ και πρέπει να φτάσει σε πολλαπλάσιο του 3: $\square=1,\,4,7$ (315, 345, 375). \quad γ) Τα δύο τελευταία ψηφία $\square 4$ πρέπει να διαιρούνται με το 4: $04, 24, 44, 64, 84$, άρα $\square=0,\,2,4,\,6,8$. \quad δ) Πρέπει να τελειώνει σε 0 (διαιρείται με 2 και 5): 170.
\sol{9} α) Λάθος: το 14 τελειώνει σε 4 και δεν διαιρείται. β) Σωστό: το 9 διαιρείται με το 3. γ) Σωστό. δ) Λάθος: το 12 διαιρείται και με τα δύο, αλλά όχι με το 24. ε) Σωστό. στ) Σωστό.
\sol{10} Το 3.465 έχει άθροισμα ψηφίων 18, άρα διαιρείται με το 3 και το 9. Τελειώνει σε 5, άρα διαιρείται με το 5. Τα δύο τελευταία ψηφία 65 δεν διαιρούνται με το 4, και δεν τελειώνει σε 0 (όχι 4, όχι 10). Σε 3 ομάδες: $3.465:3=1.155$ €. Σε 5: $693$ €. Σε 9: $385$ €.
\sol{11} Το 140 τελειώνει σε 0, άρα διαιρείται με το 2, το 5 και το 10. Τα δύο τελευταία ψηφία 40 διαιρούνται με το 4, άρα και με το 4. Το άθροισμα ψηφίων είναι 5, άρα όχι με το 3, και επομένως ούτε με το 6. Γίνεται με 2, 4, 5 και 10 στοίβες.
\sol{12} Ο αριθμός είναι $72\square 0$, με άθροισμα ψηφίων $9+\square$. Πρέπει να διαιρείται με το 9, άρα $\square=0$ ή $9$. Επειδή το τρίτο ψηφίο είναι μεγαλύτερο από το 2, είναι το 9: ο κωδικός είναι 7.290.
\sol{13} Με το 5: $60:5=12$ θέσεις. Με το 10: 6 θέσεις (10, 20, \dots, 60). Και με το 3 και με το 5: οι 15, 30, 45, 60, δηλαδή 4 θέσεις.
\sol{14} Το 2.028 είναι δίσεκτο ($28:4=7$). Το 2.000 είναι δίσεκτο (διαιρείται με το 400). Το 1.900 δεν είναι (διαιρείται με το 100 αλλά όχι με το 400). Το 2.100 δεν είναι. Το 2.026 δεν είναι (26 δεν διαιρείται με το 4).
\sol{15} Αν ένας αριθμός τελειώνει σε 0, είναι πολλαπλάσιο του 10: $10\cdot\kappa$. Επειδή $10=2\cdot 5$, έχουμε $10\kappa=2\cdot(5\kappa)=5\cdot(2\kappa)$. Άρα διαιρείται και με το 2 και με το 5.
\sol{16} Τα κομμάτια $1\cdot 999$, $4\cdot 99$ και $5\cdot 9$ διαιρούνται με το 9 και με το 3. Αρκεί να ελέγξουμε το $1+4+5+2=12$. Το 12 διαιρείται με το 3 (το 1.452 διαιρείται με το 3) αλλά όχι με το 9 (το 1.452 δεν διαιρείται με το 9).
\sol{17} Πρέπει να διαιρείται με το 4, το 9 και το 5, και τότε διαιρείται και με το 2 και με το 3. Τα πολλαπλάσια του $4\cdot 9\cdot 5=180$ είναι 180, 360, \dots Το πρώτο τετραψήφιο είναι $180\cdot 6=1.080$.
\sol{18} Το μεσαίο ψηφίο είναι το άθροισμα των δύο άκρων: $1+1=2$, $1+2=3$, $1+3=4$, $1+4=5$. Το 165 ταιριάζει ($1+5=6$) και $165=11\cdot 15$. Το 175 δεν ταιριάζει ($1+5=6\neq 7$) και $175:11$ δίνει υπόλοιπο 10.
\sol{19} Για το 3 μετράνε όλα τα ψηφία, όχι μόνο το τελευταίο. Το άθροισμα είναι $2+0+2+6=10$, που δεν διαιρείται με το 3, άρα ο 2.026 δεν διαιρείται (η διαίρεση δίνει $2.026=3\cdot 675+1$).

\lyseis{6}{Πρώτοι αριθμοί και ανάλυση: τα δομικά υλικά}
\sol{1} 12: 1, 2, 3, 4, 6, 12. \quad 18: 1, 2, 3, 6, 9, 18. \quad 30: 1, 2, 3, 5, 6, 10, 15, 30. \quad 36: 1, 2, 3, 4, 6, 9, 12, 18, 36. \quad 45: 1, 3, 5, 9, 15, 45.
\sol{2} Του 4: 0, 4, 8, 12, 16, 20. Του 7: 0, 7, 14, 21, 28, 35. Του 12: 0, 12, 24, 36, 48, 60.
\sol{3} Πρώτοι: 2, 13, 29, 41, 61, 89. Το 1 δεν είναι ούτε πρώτος ούτε σύνθετος. Σύνθετοι: 9 (3), 21 (3), 33 (3), 51 (3), 57 (3), 87 (3). Όλοι έχουν διαιρέτη το 3, γιατί το άθροισμα των ψηφίων τους (9, 3, 6, 6, 12, 15) διαιρείται με το 3.
\sol{4} 53, 59, 61, 67, 71, 73, 79.
\sol{5} α) $24=2^3\cdot 3$ \quad β) $60=2^2\cdot 3\cdot 5$ \quad γ) $72=2^3\cdot 3^2$ \quad δ) $100=2^2\cdot 5^2$ \quad ε) $150=2\cdot 3\cdot 5^2$ \quad στ) $504=2^3\cdot 3^2\cdot 7$.
\sol{6} α) $8\cdot 5=40$ \quad β) $2\cdot 9\cdot 7=126$ \quad γ) $4\cdot 3\cdot 25=300$ \quad δ) $27\cdot 11=297$.
\sol{7} α) Λάθος: ο 2 είναι άρτιος πρώτος. β) Λάθος. γ) Λάθος: $91=7\cdot 13$. δ) Σωστό. ε) Σωστό: $0=\alpha\cdot 0$. στ) Λάθος: $51=3\cdot 17$.
\sol{8} Ο 221 δεν διαιρείται με το 2, το 3 ($2+2+1=5$), το 5, το 7 ($7\cdot 31=217$) ούτε το 11 ($11\cdot 20=220$). Με το 13: $13\cdot 17=221$. Οι πρώτοι είναι 13 και 17.
\sol{9} Οι διαιρέτες του 48 είναι 1, 2, 3, 4, 6, 8, 12, 16, 24, 48: σακουλάκια με όποιον από αυτούς τους αριθμούς καραμελών. Το 47 είναι πρώτος, άρα έχει μόνο τους διαιρέτες 1 και 47: μόνο δύο τρόποι (47 σακουλάκια με 1 καραμέλα ή ένα σακουλάκι με 47).
\sol{10} Οι διαιρέτες του 60 είναι 1, 2, 3, 4, 5, 6, 10, 12, 15, 20, 30, 60. Αφαιρώ τον 1 και τον 60: μένουν 10, άρα 10 τρόποι (2, 3, 4, 5, 6, 10, 12, 15, 20 ή 30 παιδιά σε κάθε ομάδα). Το 61 είναι πρώτος, άρα κανένας τρόπος.
\sol{11} Οι πρώτοι ανάμεσα στο 20 και στο 30 είναι οι 23 και 29. Άθροισμα ψηφίων: $2+3=5$ και $2+9=11$. Ο αριθμός είναι ο 29.
\sol{12} Οι πρώτοι μέχρι το 60 είναι 17: 2, 3, 5, 7, 11, 13, 17, 19, 23, 29, 31, 37, 41, 43, 47, 53, 59. Το 2 σβήνει 29 αριθμούς (4, 6, 8, \dots, 60). Το 3 σβήνει ακόμα 9 αριθμούς: 9, 15, 21, 27, 33, 39, 45, 51, 57 (τα άλλα πολλαπλάσιά του είναι ήδη σβησμένα, γιατί είναι άρτια).
\sol{13} Για το 6: $1+2+3=6$. Για το 28: $1+2+4+7+14=28$. Σε κάθε περίπτωση το άθροισμα των διαιρετών (εκτός από τον ίδιο) ισούται με τον αριθμό.
\sol{14} $(3,\,5)$, $(5,\,7)$, $(11,\,13)$, $(17,\,19)$, $(29,\,31)$, $(41,\,43)$.
\sol{15} Οι αριθμοί είναι $1, 3, 9, 2, 6, 18, 4, 12, 36, 8, 24, 72$, δηλαδή 12 αριθμοί. Είναι όλοι οι διαιρέτες του 72.
\sol{16} Αν το 1 ήταν πρώτος, το 6 θα είχε πολλές αναλύσεις: $2\cdot 3$, $1\cdot 2\cdot 3$, $1\cdot 1\cdot 2\cdot 3$, \dots, όσες θέλεις. Η ανάλυση δεν θα ήταν μοναδική.
\sol{17} $101\cdot 103=10.100+303=10.403$. Αν δοκιμάζαμε τους πρώτους με τη σειρά (2, 3, 5, \dots, 97), θα αποτυγχάναμε 25 φορές, και μόνο ο 101 θα πετύχαινε. Με ακόμα μεγαλύτερους αριθμούς η δουλειά γίνεται τεράστια, και σε αυτό στηρίζεται η ασφάλεια των κωδικών.

\lyseis{7}{ΜΚΔ και ΕΚΠ: μοιράζω και συναντιέμαι}
\sol{1} α) Διαιρέτες του 12: 1, 2, 3, 4, 6, 12. Του 30: 1, 2, 3, 5, 6, 10, 15, 30. Κοινοί: 1, 2, 3, 6. $\MKD=6$. \quad β) Του 18: 1, 2, 3, 6, 9, 18. Του 24: 1, 2, 3, 4, 6, 8, 12, 24. $\MKD=6$. \quad γ) Του 20: 1, 2, 4, 5, 10, 20. Του 45: 1, 3, 5, 9, 15, 45. $\MKD=5$. \quad δ) Του 16: 1, 2, 4, 8, 16. Του 27: 1, 3, 9, 27. Κοινός μόνο ο 1: $\MKD=1$.
\sol{2} α) $\EKP(4,\,6)=12$ \quad β) $\EKP(6,\,15)=30$ \quad γ) $\EKP(8,\,12)=24$ \quad δ) $\EKP(5,\,7)=35$.
\sol{3} α) $24=2^3\cdot 3$, $36=2^2\cdot 3^2$: $\MKD=2^2\cdot 3=12$, $\EKP=2^3\cdot 3^2=72$. \quad β) $60=2^2\cdot 3\cdot 5$, $84=2^2\cdot 3\cdot 7$: $\MKD=2^2\cdot 3=12$, $\EKP=2^2\cdot 3\cdot 5\cdot 7=420$. \quad γ) $90=2\cdot 3^2\cdot 5$, $150=2\cdot 3\cdot 5^2$: $\MKD=2\cdot 3\cdot 5=30$, $\EKP=2\cdot 3^2\cdot 5^2=450$. \quad δ) $48=2^4\cdot 3$, $180=2^2\cdot 3^2\cdot 5$: $\MKD=2^2\cdot 3=12$, $\EKP=2^4\cdot 3^2\cdot 5=720$.
\sol{4} $12=2^2\cdot 3$, $18=2\cdot 3^2$, $30=2\cdot 3\cdot 5$. $\MKD=2\cdot 3=6$ (κοινοί παράγοντες 2 και 3, με μικρότερο εκθέτη 1). $\EKP=2^2\cdot 3^2\cdot 5=180$.
\sol{5} Πρώτοι μεταξύ τους είναι τα (α), (β), (δ) και (στ). Τα (γ) 21 και 35 έχουν κοινό διαιρέτη το 7, και τα (ε) 14 και 49 επίσης το 7.
\sol{6} α) $\MKD=7$, $\EKP=42$ \quad β) $15$ και $60$ \quad γ) $25$ και $100$. Όταν ο ένας αριθμός διαιρεί τον άλλον, το ΜΚΔ είναι ο μικρότερος και το ΕΚΠ ο μεγαλύτερος.
\sol{7} α) Λάθος: το ΜΚΔ είναι 6. β) Λάθος: το ΕΚΠ είναι 12. γ) Σωστό. δ) Σωστό. ε) Σωστό. στ) Σωστό.
\sol{8} Το ΕΚΠ(15, 20): $15=3\cdot 5$, $20=2^2\cdot 5$, άρα $2^2\cdot 3\cdot 5=60$ δευτερόλεπτα. Στα 3 λεπτά (180 δευτ.) συμπίπτουν στα 60, 120 και 180 δευτερόλεπτα: 3 φορές.
\sol{9} $84=2^2\cdot 3\cdot 7$, $126=2\cdot 3^2\cdot 7$, άρα $\MKD=2\cdot 3\cdot 7=42$. Το πολύ 42 καλάθια, με $84:42=2$ μήλα και $126:42=3$ πορτοκάλια το καθένα.
\sol{10} $180=2^2\cdot 3^2\cdot 5$, $150=2\cdot 3\cdot 5^2$, άρα $\MKD=2\cdot 3\cdot 5=30$. Μέγιστη πλευρά 30 εκ. Πλακάκια: $180:30=6$ στο μήκος και $150:30=5$ στο πλάτος, δηλαδή $6\cdot 5=30$ πλακάκια.
\sol{11} $\EKP(15,\,20)=60$ λεπτά, άρα ξανά μαζί στις 8:00. Μαζί βρίσκονται στις 7:00, 8:00, 9:00 και 10:00: 4 φορές.
\sol{12} $\EKP(4,\,6,9)$: $4=2^2$, $6=2\cdot 3$, $9=3^2$, άρα $2^2\cdot 3^2=36$. Σε 36 μέρες.
\sol{13} $(12,\,18)$: ΜΚΔ 6, ΕΚΠ 36, $6\cdot 36=216=12\cdot 18$. $(8,\,20)$: 4 και 40, $4\cdot 40=160=8\cdot 20$. $(9,\,15)$: 3 και 45, $3\cdot 45=135=9\cdot 15$. $(14,\,21)$: 7 και 42, $7\cdot 42=294=14\cdot 21$. Το γινόμενο ΜΚΔ·ΕΚΠ είναι πάντα ίσο με το γινόμενο των δύο αριθμών.
\sol{14} $84=30\cdot 2+24$, και $\MKD(84,\,30)=6=\MKD(30,\,24)$. Συνεχίζω: $30=24\cdot 1+6$ και $24=6\cdot 4+0$. Ο τελευταίος διαιρέτης πριν το υπόλοιπο 0, το 6, είναι το ΜΚΔ. Για τα 252 και 105: $252=105\cdot 2+42$, $105=42\cdot 2+21$, $42=21\cdot 2+0$. Άρα $\MKD(252,\,105)=21$.
\sol{15} $8=2^3$ και $15=3\cdot 5$ δεν έχουν κοινό παράγοντα. Το ΕΚΠ πρέπει να περιέχει όλους τους παράγοντες και των δύο, με τον μεγαλύτερο εκθέτη, και επειδή δεν υπάρχει κανένας κοινός, τους παίρνει όλους: $2^3\cdot 3\cdot 5=8\cdot 15=120$.
\sol{16} Το γινόμενο των δύο αριθμών πρέπει να είναι $6\cdot 36=216$. Ένα ζευγάρι είναι το 6 και το 36 (ΜΚΔ 6, ΕΚΠ 36). Δεύτερο ζευγάρι: 12 και 18 (ΜΚΔ 6, ΕΚΠ 36, γινόμενο 216).
\sol{17} Όχι. Το ΜΚΔ είναι μικρότερο ή ίσο από τον μικρότερο αριθμό και το ΕΚΠ μεγαλύτερο ή ίσο από τον μεγαλύτερο. Αν οι αριθμοί είναι διαφορετικοί, ο μικρότερος είναι αυστηρά μικρότερος από τον μεγαλύτερο, άρα $\MKD<\EKP$. Ισότητα έχουμε μόνο όταν οι δύο αριθμοί είναι ίσοι.

\lyseis{1}{Επανάληψη ενότητας 1}
\setcounter{chapter}{0}
\solr{1}{1} $37.002.005$. Στα εκατομμύρια: το ψηφίο των εκατοντάδων χιλιάδων είναι 0, άρα $37.000.000$.
\solr{1}{2} $4.999,\ 5.012,\ 5.021,\ 5.102,\ 10.000$.
\solr{1}{3} Δεκάδες: $86.550$ (μονάδες 9). Εκατοντάδες: $86.500$ (δεκάδες 4). Χιλιάδες: $87.000$ (εκατοντάδες 5).
\solr{1}{4} α) $25\cdot 4\cdot 17=1.700$ \quad β) $99\cdot 8=800-8=792$ \quad γ) $36\cdot(15+85)=3.600$ \quad δ) $(47+53)+(68+32)=200$.
\solr{1}{5} α) $3\cdot 36-10=98$ \quad β) $32-25=7$ \quad γ) $20-8=12$ \quad δ) $100-4\cdot 25=0$.
\solr{1}{6} $1.000=48\cdot 20+40$. Ο αριθμός είναι $15\cdot 7+6=111$.
\solr{1}{7} $100=7\cdot 14+2$: Πέμπτη και 2 μέρες, δηλαδή Σάββατο.
\solr{1}{8} Με 2 (τελειώνει σε 0), 5, 10, 3 και 9 (άθροισμα ψηφίων 9), 4 (40 διαιρείται με το 4). Όχι με το 25 (40 δεν είναι 00, 25, 50, 75).
\solr{1}{9} $5+7=12$, και χρειάζεται 18, άρα $\square=6$: $567=9\cdot 63$.
\solr{1}{10} Πρώτοι: 59 και 67. Σύνθετοι: $57=3\cdot 19$ και $63=3^2\cdot 7$.
\solr{1}{11} $360=2^3\cdot 3^2\cdot 5$. Έχει 3 διαφορετικούς πρώτους παράγοντες.
\solr{1}{12} $36=2^2\cdot 3^2$ και $60=2^2\cdot 3\cdot 5$: $\MKD=2^2\cdot 3=12$, $\EKP=2^2\cdot 3^2\cdot 5=180$.
\solr{1}{13} Ο αριθμός των στοιβών είναι ο $\MKD(36,\,60)=12$. Σε κάθε στοίβα: $36:12=3$ κάρτες με ζώα και $60:12=5$ με αυτοκίνητα.
\solr{1}{14} $\EKP(8,\,12)=24$: σε 24 μέρες.
\solr{1}{15} α) Λάθος (το 6 είναι άρτιος και δεν διαιρείται με το 4). β) Σωστό. γ) Λάθος: $91=7\cdot 13$. δ) Λάθος: $2^4=16$. ε) Λάθος: $17-(9-3)=11$, ενώ $17-9-3=5$. στ) Σωστό.
\solr{1}{16} Δίκιο έχει η Ζωή: ο πολλαπλασιασμός γίνεται πρώτος, $4\cdot 5=20$, και $6+20=26$. Για να βγει 50: $(6+4)\cdot 5=50$.
\solr{1}{17} $100=2^2\cdot 5^2$. Οι διαιρέτες είναι 1, 2, 4, 5, 10, 20, 25, 50, 100: 9 διαιρέτες. Πρώτοι είναι οι 2 και 5, δηλαδή 2.


\end{document}
